\documentclass[10.5pt, a4paper]{ctexbook} % 10.5pt = 小四字号

% 导入设定
% ===== 核心排版设置 =====
\usepackage[top=2.5cm, bottom=2.5cm, left=3cm, right=2.5cm, 
            headheight=15pt, includehead]{geometry}
\usepackage{setspace}
\usepackage{titlesec}

% ===== 视觉增强宏包 =====
\usepackage{fontspec}
\usepackage[perpage,symbol*]{footmisc}
\usepackage[dvipsnames, svgnames, x11names]{xcolor}

% ===== 内容结构宏包 =====
\usepackage{array} % for better column formatting
\usepackage{booktabs} % for nicer table rules
\usepackage[most]{tcolorbox}
\usepackage{enumitem}       % 列表定制
\usepackage{graphicx}       % 插图支持
\usepackage{amssymb}         % 示例文本

% ===== 字体配置 =====
\setmainfont{Times New Roman}
\setsansfont{Arial}
\setCJKmainfont[BoldFont=STZhongsong]{SimSun}
\setCJKsansfont{Microsoft YaHei}
\setCJKmonofont{NSimSun}

\definecolor{lessonblue}{RGB}{45, 85, 155}
\definecolor{examplebg}{RGB}{245, 248, 252}

% ===== 例句环境 =====
\tcbset{
  enhanced,
  breakable,
  sharp corners,
  boxsep=6pt,
  colback=examplebg,
  colframe=examplebg,
}

\newtcolorbox{liju}{
  breakable,
  colback=examplebg,
  colframe=examplebg,
%   borderline horizontal={1.2pt}{0pt}{lessonblue}, % horizontal rule
  before upper={%
    % \vspace{0.3\baselineskip}%
    \noindent\textbf{\color{lessonblue}例句}%
    % \vspace{0.3\baselineskip}%
    \par\noindent{\color{lessonblue}\rule{\linewidth}{1.2pt}}%
    \par\vspace{0.3\baselineskip}%
    \setlength{\parindent}{0pt}%
    \setlist[enumerate]{noitemsep, topsep=0pt, leftmargin=*}%
  },
%   after upper={\vspace{0.4\baselineskip}}
}

% 例句内容命令
\newcommand{\enzh}[2]{%
  \noindent \makebox[1em][l]{\raisebox{0.25ex}{\textcolor{lessonblue}{\large\ensuremath{\bullet}}}}%
  {#1}\par
  \noindent\hspace*{1.5em}{\small\textcolor{gray!30!black}{#2}}\par
  \vspace{0.3\baselineskip}%
}

% 正误对照命令
\newcommand{\compare}[2]{%
  % 错误句子行
  \noindent\makebox[1.8em][l]{\textcolor{red!70!black}{\large\ensuremath{\times}}}%
  {\textcolor{red!30!black}{#1}}\par
  % 正确句子行
  \noindent\makebox[1.8em][l]{\textcolor{blue!70!black}{\large\checkmark}}%
  {\textcolor{blue!30!black}{#2}}\par
  \vspace{0.3\baselineskip}% 紧凑间距（体现对照关系）
}

% 例句讲解命令
\newcommand{\explain}[3]{%
  \noindent \makebox[1em][l]{\raisebox{0.25ex}{\textcolor{lessonblue}{\large\ensuremath{\bullet}}}}%
  #1\par
  \noindent\hspace*{1.5em}{\small\textcolor{gray!30!black}{#2}}\par
  % ===== 语法讲解部分 (专业排版) =====
  \noindent\hspace*{1em}{%
    \color{lessonblue!80!black}%
    \small%
    % \raisebox{0.2ex}{\tiny\faLightbulbO}\ % 灯泡图标（可选）
    \color{lessonblue!60!black}%
    #3%
  }\par
  \vspace{0.4\baselineskip}%
}

% ==== 中文环境优化符号 ====
% ==== 音标 ====
% \newcommand{\phonetic}[1]{{\color{blue!70!black}/}{ #1}{{\color{blue!70!black}/}}}
\newfontfamily\ipafont{Charis SIL}[
  Path = fonts/,
  UprightFont = *-R,
  ItalicFont = *-I,
  BoldFont = *-B,
  BoldItalicFont = *-BI
]
\newcommand{\phonetic}[1]{{\color{blue!70!black}/\,}{\raisebox{0.25ex}{\ipafont\small\textcolor{blue!70!black}{#1}}}{{\color{blue!70!black}\,/}}}

% ==== 箭头 ====
\newcommand{\warr}{%
%   \hspace{0.15em}% 前间距
  \raisebox{-0.05ex}{\scalebox{0.9}{→}}% 提升0.2ex，缩小25%
  \hspace{0.25em}% 后间距
}

% ==== 缀字符 ====
\newcommand{\jie}[1]{%
  \raisebox{-0.1ex}{\scalebox{0.85}{–}}% 提升0.2ex，缩小15%
  \hspace{0.05em}% 微调间距
  #1%
}

% ===== 全局样式设置 =====
\ctexset{
  chapter={
    name={第,讲},
    number={\chinese{chapter}},
    format=\Huge\bfseries\color{lessonblue}\sffamily,
    aftername=\hspace{0.5em},
    beforeskip={2\baselineskip},
    afterskip={1.2\baselineskip}
  },
  section={
    format=\LARGE\bfseries\color{lessonblue!80!black}\sffamily,
    aftertitle={\hspace{0.3em}\color{lessonblue!50!black}·\hspace{0.3em}},
    beforeskip={1.2\baselineskip},
    afterskip={0.8\baselineskip}
  }
}

% 设置间距
\onehalfspacing
% \setlength{\lineskip}{24pt}
\setlength{\parskip}{6pt}
\setlength{\fboxsep}{6pt}
\setlength{\parindent}{2em}
\setlist{nosep, leftmargin=*, before=\vspace{-0.2\baselineskip}}


% 设置书名
% \titleformat{\chapter}{\zihao{-1}\bfseries}{ }{16pt}{}
% \titleformat{\section}{\zihao{-2}\bfseries}{ }{0pt}{}
\title{\zihao{0} \bfseries 英语入门}
\author{}
\date{}
% ===== 导言区结束 =====

\begin{document}
\maketitle

\frontmatter % 这行必须有 - 让前言用罗马数字页码
\chapter*{前言} % 创建无编号的"前言"标题
\addcontentsline{toc}{chapter}{前言} % 让前言出现在目录中

亲爱的同学们：

从这个学期开始，我们要学习英语了。英语是英国人的语言，随着英国在全世界的扩张，成为了全世界很多国家和地区常用的语言。很多国家的人用英语交流文化和知识，做买卖、谈生意。因此，我们有必要学好英语。

中文用汉字记录语言，英语用字母记录语言。中文的每个汉字都有自己的含义，排列组合起来表达完整的意思。而英语是由一个个单词组成的，单词的读音表达了它的意思。这与中文大不一样。中文用文字表达意思，英语用声音表达意思，字母只用来拼出单词的读音。

英语有26个字母。字母按一定的顺序排列组合起来，就成了单词。最短的单词只有一个字母，长的单词可以由几十个字母组成。单词按照一定的规则构成句子，表达完整的意思。

字母排列拼成单词，是为了记录单词的读音。我们根据汉字的拼音，就能拼出汉字的读音。但是，英语在数百年的传播和使用中经历了复杂的演变，今天，英语的拼写已经不完全反映它的读音了。不过，还是有一些大致的规则可以遵循。后续的学习中，我们会逐步掌握这些规则。

下面，就让我们开始英语的学习吧！

\vfill % 将署名推到页面底部
\begin{flushright}
编者团队 \\
\today
\end{flushright}


% ===== 正文开始 =====
\mainmatter % 关键：开始正文（阿拉伯数字页码）
\tableofcontents
% \newpage

\chapter{英语是什么？}

英语是英国人的语言，随着英国在全世界的扩张，成为了全世界很多国家和地区常用的语言。很多国家的人用英语交流文化和知识，做买卖、谈生意。因此，我们有必要学好英语。

中文用汉字记录语言，英语用字母记录语言。中文的每个汉字都有自己的含义，排列组合起来表达完整的意思。而英语是由一个个单词组成的，单词的读音表达了它的意思。这与中文大不一样。中文用文字表达意思，英语用声音表达意思，字母只用来拼出单词的读音。

英语有26个字母。字母按一定的顺序排列组合起来，就成了单词。最短的单词只有一个字母，长的单词可以由几十个字母组成。单词按照一定的规则构成句子，表达完整的意思。

字母排列拼成单词，是为了记录单词的读音。我们根据汉字的拼音，就能拼出汉字的读音。但是，英语在数百年的传播和使用中经历了复杂的演变，今天，英语的拼写已经不完全反映它的读音了。

\section{音节和音素}

英语是通过声音表达意思的。声音的组合才是英语的本体。英语是由一个个单词组成的。单词是英语里表达意思的基本单位。每个单词又是由一个或多个音节组成的。一个音节就是我们说的“发一次音”、“吐一个字”。比如，每个汉字的读音都是一个音节。音节可以说是“发音的节奏”。

汉语的一个音节，对应一个字。我们用拼音记录汉字的读音，分为声母和韵母。英语的音节也有类似的结构，但更复杂一点。英语里，一个音节由起音、音元、收音构成。其中音元是音节的核心，决定了这个音节的主要发音特征。起音和收音是音节的附属部分，起到辅助发音的作用。一个音节必然有音元，但不一定有起音和收音。

音节是音素构成的。音素是我们的嘴巴能发出的、耳朵能分辨出来的最小的声音单位。汉语里，每个声母是一个音素，每个简单的韵母也是一个音素。比如“八”的声母和韵母分别是一个音素。而复杂的韵母可以有多个音素。比如“单”的韵母由“啊”和前鼻音组成。“双”的韵母由“乌”、“啊”和后鼻音组成。英语里，一个音节的起音、音元、收音都可以由一个或多个音素构成。组成音元的音素叫做元音，组成其他部分的音素叫辅音。用来记录它们的字母，也因此分别叫做元音字母和辅音字母。

\section{英语的字母}

我们用拼音字母来拼出汉字的读音。拼音字母一共有26个。英语也用这26个字母拼出单词的读音。这26个字母叫做拉丁字母，也叫罗马字母，是古代罗马人用的字母。很多欧洲的语言都用这套字母来拼出自己的单词。

英语规定，26个字母里有5个叫做元音字母，分别是：a、e、i、o、u。其余的字母叫做辅音字母（字母y有时也用来记录元音，所以y也称为半元音字母）。英语用字母的排列组合，记录单词的读音。

英语给每个字母指定了一个读音。英语里单独读这个字母时，就用这种读法。

拉丁字母分大写和小写两种写法。英语里，两种写法都会用到。一般来说，普通的单词用小写字母表示。每句话的第一个单词，第一个字母要大写。另外，一种称为专有名词的特殊单词，第一个字母需要大写。某些情况下，比如要表示强调，或传达强烈的感情时，我们也会见到全部用大写字母拼成的单词。

\section{轻音和重音}

英语的单词由音节构成。由多个音节构成的单词，有些音节要着重去读，有些音节要轻一点读。我们称为重音节和轻音节。多音节的词里，哪些音节是重音节，哪些音节是轻音节呢？没有绝对的规律。不过，表示人和事物的名字、概念的词，往往重音节在最前。表示动作的词，往往轻音节在最前。

要注意的是，有的单词既能表示名字和概念，也能表示动作。这种单词有两种读法。读法不同，意义也不同。重音节在最前时，往往是名字和概念的词；轻音节在最前的，往往是表示动作的词。

% 动词不定式（to do）的用法：

% \begin{liju}
%   \enzh{She wants to learn English.}{她想学英语。}
%   \enzh{It's important to practice every day.}{每天练习很重要。}
%   \enzh{He promised to help me.}{他答应帮助我。}
% \end{liju}

% \section{常见错误}

% 注意以下错误用法：

% \begin{liju}
%   \compare{She wants learn English.}{She wants \underline{to} learn English.}
%   \compare{It's important practice every day.}{It's important \underline{to} practice every day.}
% \end{liju}

\section{单词的拼写}

英语用拉丁字母的排列组合记录单词。每个字母有一种或几种读法（但不一定是单独读这个字母时的读法），对应一个或几个音素。有些字母的固定组合也有固定的读法。这些常见的搭配，叫做“拼读规则”。

不过，英语的拼写也受到历史的影响。很多单词是从拉丁语、法语、希腊语等别的语言借来的，它们保留了原来的拼法，但读音却发生了变化。因此，同一个字母或组合，在不同单词里可能读法不同。

英语单词的拼写是固定的，总的来说，也是统一的。不过也有例外：极少数单词，不同国家地区的人有不同的习惯写法。

\section{音标}

英语的拼写系统历史悠久，受到多种语言影响，同一个字母或组合在不同单词中可能读法不同，比如 ea 在 head 中读 \phonetic{e}，在 seat 中读\phonetic{iː}，在 break 中又读 \phonetic{eɪ}。仅靠拉丁字母拼写，难以判断一个单词该怎么读。

为了准确地表示英语的读音，我们会使用音标来标记。相比拉丁字母拼写，音标能更准确地标识英语的读音，但音标不是用来代替拼写的，而是作为拼写的补充。我们目前采用拉丁字母拼写作为英语单词的正确记录方式。音标可以用来帮助我们准确读出英语单词，但不能作为单词来使用。就像汉语拼音帮助我们认识汉字的读音一样，音标帮助我们掌握英语单词的正确发音。我们可以借助音标学会如何朗读一个新词，但最终书写时仍然要使用传统的拉丁字母拼写。

常用的音标系统不止一种，这里我们采用一种基于国际音标系统（IPA）的记法。它能帮助我们更准确地记录单词及其音节的读音。在这个音标系统里，英语有12个元音音素，24个辅音音素。此外，还有一些按规律复合的音素，我们将其视为一个音素整体认读。合共51个音素。

\section{英音和美音}

汉语有很多方言。同一个方言里，还分不同口音。比如，北京人、天津人和南京人讲普通话，总有一些不同。英语和汉语一样，存在不同的方言和口音。其中最主要的两种口音是英国和美国的标准口音，分别称为英式英语和美式英语。这两种口音就好比英国和美国的“普通话”。这两种口音在不少单词的发音上都有所不同，但这并不影响它们都是正确、地道的英语。

使用辞典的时候，我们会发现，辞典对一个单词提供两种音标，分别是英式音标和美式音标。它们分别用来标识英式英语口音和美式英语口音的读法。

英式英语和美式英语有什么不同呢？一个最显著的标志是“卷舌音”。在美式英语中，只要单词中有字母“r”，它就几乎总是要发音，表现为一个明显的卷舌动作，音标符号为 \phonetic{r}，例如在 car（美 \phonetic{kɑː r}，英 \phonetic{kɑː}）和 teacher（美 \phonetic{ˈtiː tʃər}，英 \phonetic{ˈtiː tʃə}）中。而在英式英语中，这个 \phonetic{r} 只有在元音前才发音。另一个常见的区别体现在一些特定元音上。例如，很多在英音中读 \phonetic{ɑː} 的词，在美音中会发成 \phonetic{æ}，比如 ask、bath 和 dance。英音中的 \phonetic{ɒ}，在美音中通常发得像 \phonetic{ɑː },比如lot 的英音读\phonetic{lɒt}，美音就读\phonetic{lɑː t} 。此外，几个双元音也有不同，最典型的是“o”的发音，英音发音为 \phonetic{əʊ}的，美音则读 \phonetic{oʊ}（比如 go 的英音读\phonetic{ɡəʊ}，美音读\phonetic{ɡoʊ}），但实际听起来非常接近。

\section{开音节和闭音节}

我们已经学过单词的音节。一个音节由起音、音元、收音构成。英语里，音节又分为开音节和闭音节两类。元音结尾、没有收音的音节叫开音节，有收音的叫闭音节。

开音节是指以元音结尾的音节，没有辅音收尾。比如 go 读\phonetic{ɡoʊ}，最后一个声音是元音 \phonetic{oʊ}，听起来开阔而响亮，这就是一个典型的开音节。类似的还有 no（\phonetic{noʊ}）和 he（\phonetic{hiː}）。在开音节中，元音常常发它自己的名称音，也就是我们常说的“长音”。这种读法让人容易听清，也更容易拉长。

闭音节则是以辅音结尾的音节，声音被“关”住了。比如 cat（\phonetic{kæt}），最后的\phonetic{t} 把音节封闭起来；又如 sit（\phonetic{sɪt}）和 run（\phonetic{rʌn}），也都以辅音结束。这类音节里的元音通常读得短促有力，也就是所谓的“短音”。这是英语中非常常见的一种模式，尤其在简单词汇中广泛存在。

一种特殊的音节是元音发长音的闭音节。它的元音是长音，但后面有辅音收音。拼写时，我们往往会在辅音字母后加一个不发音的 e，用来标记这个音节是发长音的闭音节。这个规则往往适用于最后一个音节。

要注意的是，虽然单个字母或两三个字母的拼读有所谓的拼读规则，但这些规则不单一，而且我们有时难以判断一个字母是否属于字母组合、属于哪个组合，所以仅根据单词的拼写往往难以确定单词的读音。而其他的欧洲语言，比如西班牙语、意大利语，拼读规则更固定更单一，往往能通过单词的拼写确定读音。我们说西班牙语、意大利语有强拼读规则，而英语只有弱拼读规则。

\chapter{句子和单词之一}

当我们开始学习英语时，会遇到一个汉语学习中很少提及的概念——语法。语法是说一门语言时需要遵守的规则。英语的语法主要包括两个方面：首先是单词在句子中如何变化；其次是句子里不同的成分按什么顺序、以什么方式相互配合，表达完整的意思。

和中文相比，英语的语法更明确，也更重要。遵守语法是说好英语的前提。一个语法错误，往往会让句子失去意义，或者意思完全变了。

英语的语法分为词法和句法。词法规定单词在句子中如何变化；句法规定句子中的成分如何排列搭配。

和中文相比，英语的词法是全新的概念。英语在语言类型上属于“屈折语”。具体来说，英语的单词是由一串音节构成的。英语习惯通过调整音节来规定单词在句中的功能，表达特定的含义。单词的调整往往在词尾，包括增添音节、改变末音节等等。落实到书面，就是单词的字母改变或在词尾增加字母。因此，英语单词往往有多种形态。单词在句子中的形态不同，功能也不同。不同形态其实是同一个单词的不同表现。词法就是变化调整单词的法则。

举例来说，drink（喝）在句子里可以变形为drinks、drinking、drank、drunk，扮演不同的功能。这几种形态都是drink这个单词的不同表现。

和词法相比，英语的句法更好懂一些。它和中文里的规则比较像，但也有区别。汉语虽然也对语序有要求，但更注重用合适的字词来表达恰当的意思。英语句子的意思和形式、语序有更紧密的关联，用固定的语序、标记、照应来确保传达恰当的意思。汉语句子构造允许灵活宽松，英语的句子构造更要求精确严密。

因此，学习英语时，需要特别关注语法规则。英语的语法既是快速了解英语的工具，也是正确使用英语的基础。

\section{句子中的单词}

单词按顺序排列成句子。单词与单词之间留有一个字母宽的空格。和中文的句子一样，我们也把英语的句子分为四类。分别是：陈述句、疑问句、祈使句、感叹句。句子首个单词的首字母要大写，其余字母小写，除非另有要求。句子的结尾用一般以句号结束，有时也以省略号结束。疑问句以问号结束，感叹句以感叹号结束，祈使句也可以以感叹号结束。句子中可以有顿号等其他标点符号。

一个或多个句子可以用引号括起来，表示被引用。被引用的句子可以放在其他句子里作为特殊的部分，不影响整句的断句。

\begin{liju}
  \enzh{``How about tomorrow? We can go to the zoo.'' she tried to ask him out, ``I am available all day.''}{“明天怎么样？我们可以去动物园。”她想约他出去玩，“明天我一整天都有空。”}
  \enzh{``That must be Mr. Green! I know him. He is a good man.'' said the kid.}{“那位想必是格林先生了！我认识他。他是个好人。”孩子说。}
\end{liju}

如果以被引用的句子结束句子，不需要额外加引号。

句子中还可以有别的句子。我们把这样的句子称为复合句。如果句子里不含有别的句子，就称为简单句。简单陈述句是最基本的句子。我们首先学习简单陈述句，再来学习更复杂的句子。

\section{句子的成分}

首先来看简单陈述句。英语里，一句话用来介绍一个对象。这个对象叫做句子的主语，介绍的内容叫做句子的谓语。英语规定，简单陈述句只有一个主语，一个谓语，主语在谓语前面。为了好说话，我们一般把主语记作“S”，把谓语记作“P”。英语的简单陈述句的结构可以表示为“S\;+\;P”。

\begin{liju}
    \explain{Pandas are cute.}{熊猫很可爱。}{这句话的主语是“pandas”（熊猫），谓语是“are cute”（可爱）。这句话介绍熊猫，说熊猫可爱。}
    \explain{Tom went out.}{汤姆出去了。}{这句话的主语是“Tom”（汤姆），谓语是“went out”（出去了）。这句话讲述汤姆的事，说汤姆出去了。}
\end{liju}

我们还可以把句子分得更细。谓语是讲述、介绍、形容主语的内容。如果要介绍或形容主语“怎么样”，就用表语（记作“Pv”）来表示，中间用系语（记作“B”）来联系起来。句子结构为："S\;+\;B\;+\;Pv”。

\begin{liju}
    \explain{Pandas are cute.}{熊猫很可爱。}{这句话的表语是“cute”，表示形容熊猫的话。系语是“are”，把主语和表语联系起来。}
    \explain{She is happy.}{她很快乐。}{这句话的主语是“she”（她），表语是“happy”（快乐），系语是“is”。}
\end{liju}

如果要讲述或介绍主语“做了什么”，就用动语（记作“V”）来表示相关的动作。句子结构为："S\;+\;V"。

\begin{liju}
    \explain{He cried.}{他哭了。}{这句话的谓语“cried”（哭了）就是一个动语，表示主语“he”（他）做了什么。}
    \explain{Winter comes.}{冬天来了。}{这句话的谓语“comes”（来了）就是一个动语，表示主语“winter”（冬天）做了什么。}
\end{liju}


有的动语不仅说明主语做了什么，还有接受动作的对象。

\begin{liju}
    \explain{He bought a pen.}{他买了一支笔。}{这句话的谓语“bought a pen”包括动作“bought”（买）和接受动作的对象“a pen”（一支笔）。}
    \explain{I love China.}{我爱中国。}{这句话的谓语“love China”包括动作“love”（爱）和接受动作的对象“China”（中国）。}
\end{liju}

我们把接受动作的对象称为宾语（记作“O”）。以上句子里的“a pen”、“China”都是宾语。宾语总是某个动语的宾语，一般情况下接在动语后面。动语和宾语合为谓语。句子结构为：“S\;+\;V\;+\;O”。

有的动语后面有宾语，有的没有。有的时候，动语后面可以接两个宾语，用来说明动作。

\begin{liju}
    \explain{He gave me a pen.}{他给我一支笔。}{这句话的谓语包含动作“gave”（给），以及宾语“me”（我）和“a pen”（一支笔）。这是因为“给”这个动作自然需要说明“给谁”和“给什么”两个内容。}
    \explain{They called him a hero.}{他们称他为英雄。}{这句话的谓语包含动作“called”（称），以及宾语“him”（他）和“a hero”（英雄）。这是因为“称”这个动作自然需要说明“称谁”和“称为什么”两个内容。}
\end{liju}

有的时候我们还想补充说明动作的效果、程度、方法等等。这时我们会在动语后面加上一部分内容，称为补语（记为“C”）。句子结构为：“S\;+\;V\;+\;O\;+\;C”。

\begin{liju}
    \explain{He cried hard.}{他哭得很厉害。}{这句话的主语是“he”（他），谓语包括动语“cried”和补语“hard”（很厉害），说明哭的程度，表示他哭得很厉害。}
    \explain{Rain makes me sad.}{雨令我伤悲。}{这句话的主语是“rain”（雨），谓语是“makes me sad”，其中“makes”是动语，“me”是宾语，“sad”（伤悲）是补语，说明“makes”的效果：令我伤悲。}
\end{liju}

\section{英语的词类}

英语的语法把单词分为十类：名词、动词、形容词、副词、代词、冠词、介词、数词、连词、感叹词。

\begin{itemize}
    \item 名词：表示人、事物、地点等对象或概念的词，事物的名称。
    \item 动词：描述动作、变化的词。
    \item 形容词：修饰、限定名词的词。
    \item 副词：修饰、限定动词或形容词的词。
    \item 代词：指代名词，避免重复的词。
    \item 冠词：加在名词前，说明是特指还是泛指的词。
    \item 介词：介于名词前，指明其关系的词。
    \item 数词：表示数字、序数的词。
    \item 连词：连接不同内容的词。
    \item 感叹词：表达情感、叹息的词。
\end{itemize}

所有单词都至少属于十类中的一类，而有些词可以属于好几类。比如walk（走）是动词，指“走”这个动作，但也是名词，表示一次散步。还有的单词可以稍微转换形态，变成别的词类的词，比如care（关心）是动词，它可以转为形容词caring，表示“关心他人的”。我们把这种词归为一个词族。词族里的词由同一个词衍生而来。

这种分类方式贯穿在英语的思想里，不仅仅用于单词，也被用来讨论句子的成分乃至句子本身。以后的学习中，我们会见到更多的例子。

\section{名词、冠词和代词}

一门语言最基本的内容是对事物的称呼。称呼事物的词称为名词。我们用名词称呼人、事物、地点等对象和概念。到底是指所有对象、特别指定的某个对象，还是任一个对象呢？英语会给名词前加上冠词，来指出具体的范围。而对于我们已经提过的对象，我们可以用代词来指代它。

\subsection{定冠词和不定冠词}

冠词是放在名词前面的一种限定词，用来限定名词所指的范围，表明该名词是特指一个对象还是泛指一类对象。

英语的冠词有两个，分别是定冠词：the，和不定冠词：a。它们直接放在名词前面。如果名词以元音开头，那么the要读作\phonetic{ði:}，a要加上鼻音，并写为an。

定冠词表示我们谈论的是名词所指的一类对象中特定的某个或某些对象，或所有对象。这特定的对象可能是谈话者已知的，也可能是谈话不知道的，但谈话者指明了：他谈的是特定的某个或某些对象，而不是在泛泛而谈。

不定冠词则表示泛泛而谈：谈论的是名词所指的一类对象中某个不特定的对象。谈话者并不关心具体是哪个对象。

名词前面不一定要加冠词。不加冠词也称为零冠词。这时的名词通常表示普遍、广泛的概念，或者表示独一无二的概念。

\subsection{普通名词和专有名词}

我们用名词称呼对象和概念，根据对象的种类，又可以分为专有名词和普通名词。专有名词是为某个对象特别指定的名称，它专属于这个对象。其他的对象可能和它同属一类，但并不用这个名称称呼。普通名词则用来指一类对象，这类对象里面的每一个都是这个名称。

专有名词的首字母常常需要大写，而普通名词的首字母通常小写。但专业名词也有首字母小写的，比如。要注意的是，同样的单词，首字母大写后变为专有名词，可能表示完全不一样的意思。比如：polish表示“抛光、擦亮”，而Polish表示“波兰人、波兰语”。

举例来说，人名、国家名、地点名、机构名、品牌名、建筑名都是专有名词。当我们为一个对象起了专门的名称，它就是独一无二的。当然，这个准则是为了说明一般情况下普通名词和专有名词的判定方式。在一些特殊的场合里，专有名词可能对应两个对象。比如，侦探小说或科幻小说里，可能出现两个一模一样的人（比如克隆人），都用同一个名字称呼。但这时候的名字同样是专有名词。

\subsection{可数名词与不可数名词}

名词又分为可数名词和不可数名词。英语把事物的观念分为两类，一类概念所指的对象可以计数的个体，称为可数的。对应的名词称为可数名词。另一类概念所指的对象没有可以计数的个体，称为不可数的。对应的名词称为不可数名词。

不可数名词大致有几类：
\begin{enumerate}
    \item 直观上无法分出最小单元的存在：water, milk, coffee, air, gas, gold, iron, wood, plastic, cotton, wool, silk, paper, cement, homework, music, time
    \item 细微颗粒构成的物质，难以计数的：sugar, rice, dust, flour, salt, wheat
    \item 抽象概念：happiness, freedom, knowledge, patience, courage, mathmatics, physics
    \item 集合概念：furniture, hair, beard, lugguage, equipment, traffic, money, jewelry, staff
\end{enumerate}

有些名词既有作为可数名词的含义，也有作为不可数名词的含义。比如chicken指鸡肉时是不可数名词，指小鸡时是可数名词。

\subsection{单数与复数}

可数名词区分单数和复数。表示一个对象的时候用单数，表示超过一个对象的时候用复数。

\subsection{集体名词和个体名词}

集体名词指代作为一个整体的一组人或事物。它本身是可数的名词。强调整体时，一个集体视为单数；强调内部的成员、成分时，一个集体也视为复数。比如：

\begin{liju}
    \explain{The team is winning. }{队伍正领先。}{强调“the team”作为一个整体，用单数。}
    \explain{The team are arguing among themselves. }{队内在讨论。}{强调“the team”的内部成员这个集体，用复数。}
\end{liju}

\textbf{常见集体名词}

\begin{itemize}
    \item people, police, cattle, the + 形容词（表示一类人）：总是复数，不加-s，前面不能加a或数词。
    \item family, class, staff, audience, crew, government, committee：可单可复。
\end{itemize}

people作为名词既有作为集体名词的含义：人民，也有作为个体名词的含义：民族。

\section{冠词的用法}

不同的名词，在不同的环境下，需要添加合适的冠词（或不添加冠词），来准确表达意思。

当我们讨论可数普通名词时，如果讨论的是作为类别的整体和概念，那么用“the+单数”；如果讨论的是作为群体的整体，那么不添加冠词并用复数；如果讨论特定的某个或某些，就用“the+单/复数”；而如果只谈论其中任意一个，就用“a/an+单数”；如果要介绍或谈论一类对象，选一个出来作为例子，但不需要特别指定，这时就用“a/an+单数”。

\begin{liju}
    \explain{The panda eats bamboo.}{熊猫吃竹子。}{表示作为整体的“熊猫”这种动物，用“the+单数”。}
    \explain{Pandas are cute.}{熊猫很可爱。}{表示作为群体的“熊猫”，不加冠词的复数。}
    \explain{I saw a panda on the tree.}{我看见树上有一只熊猫。}{表示不特定的一只熊猫，用，用“a+单数”。}
    \explain{A panda can spend more than 10 hours a day on eating.}{熊猫每天能花10个小时吃东西。}{用一只熊猫来说明熊猫的习性，用“a+单数”。}
    \explain{The panda on the tree went down.}{树上的熊猫下来了。}{表示特定的那只熊猫，用“the+单数”。}
    \explain{The pandas in the zoo seem happy.}{动物园里的熊猫看起来很快乐。}{表示特定的多只熊猫，用“the+复数”。}
\end{liju}

专有名词前一般不加冠词。

\begin{liju}
    \enzh{Beijing is the capital of China.}{北京是中国的首都。}
    \enzh{Knowledge is power.}{知识就是力量。}
    \enzh{Jupiter is the largest planet in the solar system.}{木星是太阳系里最大的行星。}
\end{liju}

江河、海洋、山脉、群岛、景点、名建筑等专有名词前通常需要加定冠词。

\begin{liju}
    \enzh{The Pacific Ocean is the biggest ocean in the world.}{太平洋是世界上最大的海洋。}
    \enzh{The Eiffel Tower is the most visited tourist attraction in the world.}{埃菲尔铁塔是世界上最多人访问的旅游景点。}
    % \enzh{The Rocky Mountains contain many national parks.}{落基山脉里有许多国家公园。}
\end{liju}

表示同姓名的几个人中的一个或特定某个，要加冠词。“the+人名复数”表示一家人。

\begin{liju}
    \enzh{A Mr. Smith just called, saying he needs to meet you now.}{有位史密斯先生打电话来，说他现在就要见您。}
    \enzh{The Smith you're looking for no longer lives here.}{你要找的那位史密斯已经不在这里住了。}
    \enzh{When the Smiths showed up, he found himself cornered.}{史密斯一家出现时，他知道自己逃不掉了。}
\end{liju}

日、月、地球前用定冠词：

\begin{liju}
    \enzh{The sun is shining.}{阳光灿烂。}
    \enzh{The moon rose above the horizon.}{月亮上来了。}
    \explain{The Earth is our home.}{地球是我们的家园。}{这里的Earth并不作为宇宙中的天体，仅仅表示地球本身。}
    \explain{Earth orbits the Sun.}{地球绕着太阳转。}{这里涉及到别的天体，故Earth视为天体，前面不加冠词。}
\end{liju}


描述性的名称或复数名称作为专有名词，加定冠词。

\begin{liju}
    \enzh{New York is a city in the United States, not the United Kingdom.}{纽约是美国的城市，不是英国的城市。}
\end{liju}

仅有的几个例外包括the Vantican, the Sahara。它们在历史上作为普通名词使用。

不可数普通名词前不能加不定冠词，只能加the或不加冠词。

特指谈话者都知道的特定对象，加定冠词。

\begin{liju}
    \explain{Would you pass me the salt please?}{请把盐拿给我。}{双方都知道要拿的盐在哪里。如果对方不知道去哪里拿，则不能这么说。}
    \explain{Please pass me the salt on the table.}{请把桌上的盐拿给我。}{指定了是哪里的盐，因此用定冠词。}
    \explain{The water in the pool is too cold.}{池子里的水太冷了。}{双方都知道指的是哪个池子里的水，因此用定冠词。}
\end{liju}
 
指作为类别的一般概念时不加冠词。

\begin{liju}
    \enzh{Salt is a stable substance.}{盐是一种稳定的物质。}
    \enzh{Poverty is rampant here.}{贫穷在此肆虐。}
\end{liju}

\section{名词的复数规则}

一般在词尾加上\phonetic{s}或\phonetic{z}音，表示复数。
\begin{enumerate}
    \item 在\phonetic{p}，\phonetic{k}，\phonetic{t}，\phonetic{f}，\phonetic{θ}等辅音之后读作\phonetic{s}（少数情况\phonetic{f}变成\phonetic{v}）
    \item 在\phonetic{b}，\phonetic{d}，\phonetic{g}，\phonetic{n}，\phonetic{l}，\phonetic{m}，\phonetic{n}，\phonetic{ð}，\phonetic{r}等辅音以及元音之后读作\phonetic{z}
    \item 在\phonetic{s}，\phonetic{z}，\phonetic{ʃ}，\phonetic{tʃ}，\phonetic{ʒ}，\phonetic{dʒ}等辅音之后读作\phonetic{ɪz}或\phonetic{əz}
    \item 不规则名词：单列说明（man、foot、mouse，单复同形等等）
\end{enumerate}

在拼写上表现为：
\begin{enumerate}
    \item 大多数情况下，在词末直接加上字母s，表示\phonetic{s}或\phonetic{z}音
    \item 以\jie{s}，\jie{x}，\jie{ch}，\jie{sh}结尾的名词加\jie{es}，读作\phonetic{ɪz}或\phonetic{əz}
    \item 一些以\jie{f}或\jie{fe}结尾的名词要换成\jie{ves}，读作\phonetic{vz}（还有一些可换可不换）
    \item "辅音字母+\;y"结尾的名词，改y为i再加\jie{es}，读作\phonetic{z}
    \item “辅音字母+\;o”结尾的名词，加\jie{es}，读作\phonetic{z}，但有一些词只加\jie{s}，还有少数“元音字母+\;o”结尾的名词也可以加\jie{es}
    \item 不规则名词：单列说明（man、foot、mouse，单复同形等等）
\end{enumerate}


\textbf{常见的复数不规则的名词}
\begin{enumerate}
    \item man \warr men，woman \warr women，Frenchman \warr Frenchmen, fireman \warr firemen
    \item foot \warr feet, tooth \warr teeth, goose \warr geese
    \item mouse \warr mice, louse \warr lice
    \item ox \warr oxen, child \warr children
\end{enumerate}

\textbf{单复同形的词}
\begin{enumerate}
    \item sheep, deer, reindeer, buffalo, antelope, bison, fish, crab, cod, trout
    \item kin, offspring
    \item craft, spacecraft, aircraft
    \item series, means, species, headquarters, barracks, crossroads, gallows
    \item Chinese, Cantonese, Burmese, Japanese, Swiss
\end{enumerate}

\textbf{拉丁语借词的复数规则}
\begin{enumerate}
    \item \jie{a}结尾的词，加\jie{e}：formula \warr formulae，antenna \warr antennae
    \item \jie{ex}或\jie{ix}结尾的词，变为\jie{ices}：index \warr indices, apex \warr apices, appendix \warr appendices, matrix \warr matrices, vortex \warr vortices
    \item \jie{us}结尾的词，变为\jie{i}：alumnus \warr alumni, stimulus \warr stimuli, cactus \warr cacti, genius \warr genii, nucleus \warr nuclei, focus \warr foci, radius \warr radii, syllabus \warr syllabi, terminus \warr termini
    \item \jie{us}结尾的词，少数变为\jie{era}或\jie{ora}：genus \warr genera, corpus \warr corpora, tempus \warr tempora, opus \warr opera
\end{enumerate}

\textbf{希腊语借词的复数规则}
\begin{enumerate}
    \item 以\jie{is}结尾的词，变为\jie{es}：axis \warr axes, analysis \warr analyses, thesis \warr theses, hypothesis \warr hypotheses, crisis \warr crises, basis \warr bases, diagnosis \warr diagnoses, oasis \warr oases, synopsis \warr synopses, emphasis \warr emphases
\item 以\jie{on}或\jie{um}结尾的词，变为\jie{a}：datum \warr data, bacterium \warr bacteria, criterion \warr criteria, medium \warr media, phenomenon \warr phenomena, memorandum \warr memoranda, erratum \warr errata, curriculum \warr curricula
\end{enumerate}

意大利语、希伯来语等外来语借词也按其固有方式变为复数。

日常使用时，由于大家不熟悉拉丁语、希腊语借词的变法，也会按普通英语规则变为复数，比如appendix \warr appendixes, formula \warr formulas。学术和专业环境中通常仍然按拉丁语、希腊语规则变为复数。

\section{代词的用法}

代词是用来指代人和事物的词。代词主要用来避免重复提及同一个东西。代词有很多种，不同的代词用于代指不同的对象。

\begin{itemize}
    \item 人称代词：I、we、you、he、she、it、they、me、us、you、him、her、it、them
    \item 所属代词：my、our、your、his、her、its、their、mine、ours、yours、his、hers、its、theirs
    \item 反身代词：myself、ourselves、yourself、himself、herself、itself、themselves、oneself
    \item 相互代词：each other 、one another
    \item 指示代词：this、that、these、those、such、same
    \item 疑问代词：who、whom、whose、which、what
    \item 关系代词：who、whom、whose、which、that
    \item 不定代词：some、any、none、all、one、every、many、someone、something、anything
\end{itemize}

\section{你我他}

用来代指人的代词是人称代词。中文的人称代词是你我他（她它），如果代指多人，就加“们”。英语的人称代词比中文更复杂一些。

首先给出人称的分类。“我”和“我们”叫做第一人称；“你”和“你们”叫做第二人称；“他（她它）”和“他（她它）们”叫做第三人称。

代词出现在主语，表示你我他（她它）。

\begin{liju}
    \enzh{I am Chinese.}{我是中国人。}
    \enzh{You are Chinese.}{你是中国人。}
    \enzh{We are Chinese.}{我们是中国人。}
    \enzh{You are Chinese.}{你们是中国人。}
    \enzh{He is Chinese.}{他是中国人。}
    \enzh{She is Chinese.}{她是中国人。}
    \enzh{They are Chinese.}{他们是中国人。}
    \enzh{It is a dog.}{它是一只狗。}
    \enzh{They are hand-made.}{它们是手工制品。}
\end{liju}

代词出现在宾语，表示你我他（她它）。

\begin{liju}
    \enzh{He likes me.}{他喜欢我。}
    \enzh{She likes you.}{她喜欢你。}
    \enzh{He saved us.}{他救了我们。}
    \enzh{I love you.}{我爱你们。}
    \enzh{I followed him.}{我跟踪他。}
    \enzh{I know her.}{我认识她。}
    \enzh{We joined them.}{我们加入了他们。}
    \enzh{He ate it.}{他吃了它。}
    \enzh{I threw them away.}{我把它们扔掉了。}
\end{liju}

\section{主格和宾格}

英语的人称代词有主格和宾格，用来标记代词用作主语还是宾语。比如，我们听到us的时候，就知道它用作宾语。不过，不是每个代词的主格和宾格都不同。比如第二人称单数的主格和宾格都是you。古代英语中，第二人称单数的主格和宾格分别是thou和thee，但今天我们已经不用了。

\section{这和那}

指代事物的代词叫指示代词。中文的指示代词是“这”和“那”，指代多个就加“些”。英语的指示代词是this（这）和that（那），指代复数用these和those。

\begin{liju}
    \enzh{This is a pen.}{这是一支笔。}
    \enzh{That is a duck.}{那是一只鸭。}
    \enzh{These are cars.}{这些是车。}
    \enzh{Those are clouds.}{那些是云。}
    \enzh{I choose this.}{我选这个。}
    \enzh{I need these.}{我需要这些。}
    \enzh{He has that.}{他拥有那个。}
    \enzh{She wants those.}{她想要那些。}
\end{liju}

\chapter{动词与时态}

动词是英语的关键成分。谓语的主要部分就是动词。动词的形态和屈折变化承载了大量的语法功能。因此，按照语法正确使用动词，是说好英语的一大关键。

一般来说，一个简单句子里只有一个谓语，谓语里只有一个核心动作，我们称对应的动词为句子的主动词。句子中其他成分里也许会有不同形态的动词，但主动词是全句的核心要点。解析英语的句子，应该以寻找主动词为先。确定了主动词，就能以它为线索，理解整个句子。

\section{动词的种类}

英语的语法把动词大致分为四类：普通动词、系动词、助动词和情态动词。大部分动词都是普通动词，而系动词、助动词和情态动词是一些含义基本、用途广泛，但又在语法上比较特殊的词。

\subsection{系动词}

中文里的系动词就是“是”。它用来主语的归属、状态等等。狭义来说，英语的系动词只有一个，叫做“be”，因此也会被叫做“be动词”。但它最常见的形态并不是“be”，而是is、are、am等等。这是因为它作为谓语，要配合主语而变化，这种现象叫做主谓搭配。

\begin{liju}
    \enzh{I am a boy.}{我是男孩。}
    \enzh{You are a girl.}{你是女孩。}
    \enzh{We are happy.}{我们很快乐。}
    \enzh{You are happy.}{你们很快乐。}
    \enzh{He is a boy.}{他是男孩。}
    \enzh{She is a girl.}{她是女孩。}
    \enzh{They are happy.}{他们很快乐。}
    \enzh{It is a box.}{它是个盒子。}
    \enzh{They are new.}{它们是新的。}
\end{liju}

以上是某个常见情况下系动词的主谓搭配。我们以后还会看到，be有更多的形态，承载更多的意义。

很多时候，我们还会把一些起连接功能的动词也叫做系动词。它们的特征是后接形容词，表示主语的状态、属性。比如：

\begin{liju}
    \enzh{He looks nice.}{他看起来人不错。}
    \enzh{The soup tastes salty.}{汤尝起来很咸。}
    \enzh{The flower smells good.}{花闻起来很香。}
    \enzh{I become confident.}{我变自信了。}
    \enzh{The sky turns gray.}{天灰了。}
    \enzh{The apple stays fresh.}{苹果还新鲜。}
\end{liju}

不过，从动词形态变化的规律来说，它们与be动词并不相同，和普通动词一致。

\subsection{助动词}

助动词是英语语法的特点，中文里没有助动词的概念。它们辅助主动词完成形态变化，实现语法功能，本身并无意义。

英语的助动词有四个：be、do、have、will。它们往往在主动词前出现，辅助它构成完整的形态变化。

要注意的是：be、do和have本身也是普通动词，有实在意义。所以要注意区分句子里的do和have到底是有实际意义的动词，还是别的动词的助动词。

\subsection{情态动词}

情态动词和助动词很像，也是辅助主动词的词。它们和助动词的区别是，它们并不是用来实现语法功能，而是有自身的意义。中文也有类似的动词。我们来看一些例子：

\begin{liju}
    \enzh{He can swim.}{他会游泳。}
    \enzh{I can help you.}{我可以帮你。}
    \enzh{I will help you.}{我会帮你。}
    \enzh{They should go to school.}{他们应该上学。}
    \enzh{He must go now.}{他得走了。}
\end{liju}

可以看出，这一类动词用来表示一种态度、意愿、情势，因此叫情态动词。

英语中常见的情态动词有：can、must、may、will、shall、dare。它们的特点包括：
\begin{enumerate}
    \item 情态动词后直接主动词。
    \item 无法单独做主动词。
\end{enumerate}

不过也有例外。dare本身也可以做普通动词，有实在意义。

\section{人称与动词}

主语是第三人称单数时，动词后要加\phonetic{s}或\phonetic{z}音，规则同名词复数。如果动词以\phonetic{s}或\phonetic{z}音结尾，则加\phonetic{ɪz}或\phonetic{əz}。这个现象称为动词的变位。

在拼写上表现为：
\begin{enumerate}
    \item 大多数情况下，在词末直接加上字母s，表示\phonetic{s}或\phonetic{z}音
    \item 以\jie{s}，\jie{x}，\jie{ch}，\jie{sh}结尾的动词加\jie{es}，读作\phonetic{ɪz}或\phonetic{əz}
    \item “辅音字母+\;y”结尾的名词，改y为i再加\jie{es}，读作\phonetic{z}
    \item “辅音字母+\;o”结尾的名词，加\jie{es}，读作\phonetic{z}
    \item 不规则名词：单列说明（be、 do、have）
\end{enumerate}

\subsection{不规则变位}

有一些动词的变位是不规则的。

\begin{enumerate}
    \item  do \warr does （“o”的读音变了）
    \item  have \warr has
    \item  be的情况最复杂，一般情况的第三人称单数是is
\end{enumerate}

\section{动词和否定句}

简单陈述句可以表达肯定的意思，也可以表达否定的意思。英语中有固定的方法来表达否定的意思。具体来说，英语通过在动词前后加上否定词not来表达否定的意思。

如果只有be动词，那么通过在be后面加not表达否定。

\begin{liju}
    \enzh{I am happy. \warr I am not happy.}{我很高兴。 \warr 我不高兴。}
    \enzh{He is fine. \warr He is not fine.}{他很好。\warr 他不好。}
    \enzh{She is tall. \warr She is not tall.}{她很高。\warr 她不高。}
\end{liju}

如果有助动词，那么在第一个助动词后加not表达否定。

\begin{liju}
    \enzh{He has finished the test. \warr He has not finished the test.}{他已经完成测试了。 \warr 他并未完成测试。}
    \enzh{He will be happy. \warr He will not be happy.}{他会快乐的。 \warr 他不会快乐的。}
    \enzh{He will be watching the game. \warr He will not be watching the game.}{他会一直关注比赛的。 \warr 他不会一直关注比赛的。}
    \enzh{I have been sleeping well. \warr I have not been sleeping well.}{我最近睡得不错。 \warr 他们我最近睡得不好。}
\end{liju}

如果有情态动词，那么在情态动词后加not表达否定。

\begin{liju}
    \enzh{He can swim. \warr He can not swim.}{他会游泳。 \warr 他不会游泳。}
    \enzh{You must go. \warr You must not go.}{你必须走。 \warr 你不能走。}
    \enzh{They should come here earlier. \warr They should not come here earlier.}{他们应该早点到。 \warr 他们不应该早点到。}
\end{liju}

要注意的是：这是语法形式上的否定，但并不是完全在意义上相互矛盾的否定（逻辑上的矛盾关系）。比如：“You must go”表达强烈的意愿、建议（你必须要走），强调必要性。和它意义相互矛盾的否定应该是否定这种必要性，即“你没必要走”。形式上的否定更接近对意愿内容的反对，即赞同必要性，但反转意愿的内容：“You must not go”表达同样强烈的意愿、建议，但内容是相反的，即“你必须不走/留下来”。如果要否定必要性而保持意愿的内容不变，则应该同时把情态动词换掉，改为“You need not go”（你不必走）。不过这种用法现在已经很罕见，人们更多用“You do not have to go”或“You do not need to go”来表示“你不必走”。关于如何对情态动词句子进行逻辑上的否定，我们在后面再详细介绍。

如果只有普通动词，那么通过在动词前加助动词do，并把主动词的形态变化转移到do上，把主动词改回原形，最后在形态变化后的do后面加not表示否定。

\begin{liju}
    \enzh{I like rabbits. \warr I do not like rabbits.}{我喜欢兔子。\warr 我不喜欢兔子。}
    \enzh{He makes breads. \warr He does not make breads.}{他做面包。 \warr 他不做面包。}
    \enzh{He cried. \warr He did not cry.}{他哭了。 \warr 他没哭。}
    \enzh{They lost the game. \warr They did not lose the game.}{他们输掉了比赛。 \warr 他们没输掉比赛。}
\end{liju}

\section{动词和疑问句}

前面说的句子都是陈述句。英语的疑问句是怎么样的呢？

英语的疑问句总以问号结尾。英语把疑问句分为两种：一般疑问句和特殊疑问句。一般疑问句指用“是”或“不是”来回答的疑问，比如“你今天上学了吗？”“太阳是圆的吗？”等。特殊疑问句指不用“是”或“不是”来回答的疑问，比如“你叫什么名字？”“这个多少钱？”“今天星期几？”“哪个比较好吃？”等。

陈述句的语序是“S+P”，也就是主语加谓语。如果谓语里有表语、宾语或补语，那么放在动词后面。疑问句则需要把这个语序变换一下，我们称之为倒装。疑问句是一种倒装句。

\subsection{一般疑问句的变化规则}

一般疑问句会把谓语动词放在句首。如果谓语动词是be动词，那么be动词要放在句首，其他词位置不变。

\begin{liju}
    \explain{I am a student. \warr Am I a student?}{我是学生。 \warr 我是学生吗？}{把am移到句首。}
    \explain{He is a boy. \warr Is he a boy?}{他是男孩。 \warr 他是男孩吗？}{把is移到句首。}
    \explain{They are Chinese. \warr Are they Chinese?}{他们是中国人。 \warr 他们是中国人吗？}{把are移到句首。}
\end{liju}

如果谓语动词是助动词加普通动词，那么把助动词放到句首，其他词位置和顺序不变。

\begin{liju}
    \explain{He has finished the test. \warr Has he finished the test?}{他完成测试了。 \warr 他完成测试了吗？}{把助动词has移到句首。}
    \explain{They will come to see you. \warr Will they come to see me?}{他们明天会来看你。 \warr 他们明天会来看我吗？}{把助动词will移到句首。}
\end{liju}

如果谓语动词是情态动词加普通动词，那么把情态动词放到句首，其他词位置和顺序不变。

\begin{liju}
    \explain{He can swim. \warr Can he swim?}{他会游泳。 \warr 他会游泳吗？}{把情态动词can移到句首。}
    \explain{You shall sleep. \warr Shall you sleep?}{你该睡觉了。 \warr 你会睡觉吗？}{把情态动词may移到句首。}
    \explain{He must finish the test. \warr Must he finish the test?}{他必须完成测试。 \warr 他必须完成测试吗？}{把情态动词must移到句首。}
\end{liju}

% 要注意的是，这种变换方法只是语法上的。如果要让一般疑问句表达对应于陈述句的意思，有些情况下需要换用别的情态动词。具体的方法和规律，我们在后面的章节再详细介绍。

普通动词的一般疑问句，要在句首增加一个do作为助动词，其他词语位置和顺序不变。要注意的是，原本普通动词在陈述句中的形态，会转移到句首的do上，而普通动词要变回原形。具体如何变化，我们将在后面详细介绍。

\subsection{特殊疑问句的变化规则}

特殊疑问句有很多种，取决于对什么进行提问。语法上来说，取决于提问的对象在原本的陈述句里作什么成分。这里只介绍针对主语、宾语和谓语动词提问的特殊疑问句。

针对宾语的特殊疑问句，需要在一般疑问句的基础上，用疑问代词代替宾语，然后将它放到句首。针对人的疑问词用who，针对物（人以外）用what。比如，如果针对宾语提问：

\begin{liju}
    \explain{You are Tom. \warr Are you Tom? \warr Are you who? \warr Who are you?}{你是汤姆。 \warr 你是谁？}{针对宾语Tom提问，因此把Tom替换为who，然后放到句首。}
    \explain{This is a pen. \warr Is this a pen? \warr Is this what? \warr What is this?}{这是一支笔。\warr 这是什么？}{针对宾语a pen提问，因此把a pen替换为what，然后放到句首。}
    \explain{You can see a cat. \warr Can you see a cat? \warr Can you see what? \warr What can you see?}{你可以看到一只猫。 \warr 你能看到什么？}{针对宾语a cat提问，因此把a cat替换为what，然后放到句首。}
    \explain{The car hits a tree. \warr Does the car hit a tree? \warr Does the car hit what? \warr What does the car hit?}{汽车撞了一棵树。 \warr 汽车撞了什么？}{}
\end{liju}

如果针对主语提问，则在陈述句的基础上，直接用疑问代词代替主语：

\begin{liju}
    \explain{Tom is good at maths. \warr Who is good at maths?}{汤姆擅长数学。 \warr 谁擅长数学？}{针对主语Tom提问，因此把Tom替换为who。}
    \explain{The car hits a tree. \warr What hits a tree?}{汽车撞了一棵树。 \warr 什么撞了一棵树？}{针对主语the car进行提问，因此把the car替换为what。}
    \explain{I can open this door. \warr Who can open this door?}{我能开这扇门。 \warr 谁能开这扇门？}{针对主语 I 进行提问，因此把 I 替换为who。}
\end{liju}

英语语法中，特殊疑问句的疑问词要么针对名词（或代词），要么针对补语，因此，针对谓语动词进行提问，只能用比较简略的方法，也就是问“他干了什么”这样的问题。比如：

\begin{liju}
    \enzh{He has finished the test. \warr What has he done?}{他会游泳。 \warr 他会干什么？}
    \enzh{He can swim. \warr What can he do?}{他会游泳。 \warr 他会干什么？}
    \enzh{Tom teaches English in Beijing. \warr What does Tom do?}{汤姆在北京教英语。 \warr 汤姆做什么？}
\end{liju}

可以看到，一般来说，针对谓语进行提问，就是把谓语的主动词换成do（形态随原来的动词），有助动词和情态动词的话则保留，其他部分省略。如果谓语有助动词或情态动词，就把助动词、情态动词放到主语前。如果没有，就在句首增加一个do作为助动词，并把形态变化转移到这个do上。最后在句首添加what作为疑问词。

针对句子的其他成分提问的疑问句，规则比较复杂，涉及到更多的概念。我们会在后面的章节里继续介绍。

\section{动词的时态}

英语动词最重要的特征是时态。时态是印欧语言中广泛存在的概念，但中文里并没有时态。时态体现了西方文化对自然、时空的认知方法。因此，我们要暂时放弃“在中文里找类似的东西”的想法，从零开始，理解时态。

英语里一共有16个时态，由“四时”和“四态”交织组成。其中10个是较常用的，其余6个已经比较罕见，或者在某些语法现象中配合出现（表格中以▲标记）。我们会从最简单最常见的时态开始，逐步来了解、掌握它们。

\vspace{8pt}
\begin{tabular}{|c|c|c|c|c|}
\hline
\textbf{时态} & \textbf{现在时} & \textbf{过去时} & \textbf{将来时} & \textbf{过去将来时} \\
\hline
\textbf{一般态} & 一般现在时 & 一般过去时 & 一般将来时 & 一般过去将来时 \\
\hline
\textbf{进行态} & 现在进行时 & 过去进行时 & 将来进行时 & ▲ 过去将来进行时 \\
\hline
\textbf{完成态} & 现在完成时 & 过去完成时 & ▲ 将来完成时 & ▲ 过去将来完成时 \\
\hline
\textbf{完成进行态} & 现在完成进行时 & ▲ 过去完成进行时 & ▲ 将来完成进行时 & ▲ 过去将来完成进行时 \\
\hline
\end{tabular}

英语动词时态的“四时”是：现在时、过去时、将来时和过去将来时。它们和时间有关，但并不单纯表示现在、过去和将来，而是强调与主语的时间关联。配合“四态”成为时态后，可以表现动作和对象在时间和状态上的联系。不带时态的动词称为动词原形。带有时态的动词，或者自身发生屈折变化，或者由助动词辅助表达时态。

作为表示核心动作的成分，句子的谓语动词总带有时态。普通动词、系动词单独出现时，自身带有时态，或由助动词表示时态。情态动词辅助主动词出现时，情态动词带有时态，而主动词不带时态，保留原型。在其他成分中，动词可以不带时态。

\section{一般现在时}

下面我们来学习最简单最常见的时态：一般现在时。一般现在时的情态动词和动词原形一样。一般现在时的普通动词，大多数时候形态和动词原形一样，仅在主语是第三人称单数时，需要稍作变换。具体规则已经介绍过了，主要是在词尾加\jie{s}或\jie{es}。对于be动词，随着人称变化，我们有以下的变位规律：

\begin{liju}
    \enzh{I am a boy.}{我是男孩。}
    \enzh{You are a girl.}{你是女孩。}
    \enzh{We are happy.}{我们很快乐。}
    \enzh{You are happy.}{你们很快乐。}
    \enzh{He is a boy.}{他是男孩。}
    \enzh{She is a girl.}{她是女孩。}
    \enzh{They are happy.}{他们很快乐。}
    \enzh{It is a box.}{它是个盒子。}
    \enzh{They are new.}{它们是新的。}
\end{liju}

什么场合下用一般现在时呢？最根本的用法是：描述刚刚发生的事，用一般现在时。也就是说，事情刚刚发生，你看着事情发生，然后说这件事发生了，那么可以用一般现在时。

\begin{liju}
    \enzh{``It rains.'' said the little girl.}{“下雨了。”小女孩说。}
    \enzh{``He shoots! He scores!''}{“他起脚射门！他得分了！”}
\end{liju}

然而，这种用法如今已经很罕见了，一般只在实时直播时使用，或者在剧本、文学作品叙事中使用。

常见的用法是：描述“总是这样”的事。也就是说，描述常识、真理、习惯、常见的动作或存续状态时，应该用一般现在时。

\begin{liju}
    \explain{Two plus two equals four.}{二加二等于四。}{这是一个数学上的真理，因此用一般现在时。}
    \explain{Tigers eat meat.}{老虎吃肉。}{这是关于动物的常识，因此用一般现在时。}
    \explain{Light travels faster than sound.}{光跑得比声音快。}{这是自然规律，因此用一般现在时。}
    \explain{We go shopping every weekend.}{我们每周末都去购物。}{这是在说一种习惯，因此用一般现在时。}
    \explain{He practices a lot.}{他经常练习。}{这是在说一个观察到的常见现象，因此用一般现在时。}
\end{liju}

这里需要注意动词本身的性质。我们可以把动词分为两类：点动词和存续动词。点动词表示动作的发生是一个时间点，比如knock（敲）、go（去）、throw（抛）等。存续动词表示动作的发生是存续的一段时间，通常表示状态，比如like（喜欢）、know（知道）、stay（停留）、live（居住）。要注意的是，在点动词和存续动词的分类上，英语和中文的动词不完全对等。

对于点动词来说，使用一般现在时，表示这个动作经常发生，是某种习惯；或者形容某种规律、真理。对于存续动词来说，使用一般现在时，表示事物处于这个状态，并可能一直如此。

要注意的是，这种用法可以用于客观描述，也可以用来表达一种主观的态度和观点，就是说话的人认为“总是这样”，认为这种经常发生的事，乃至是常识、真理。至于事实上是否如此，则不一定。

\begin{liju}
    \explain{You like him.}{你喜欢他。}{这里用一般现在时，表达一种判断，即认为“你喜欢他”是真的。}
    \explain{He knows.}{他知道了。}{这里用一般现在时，表达一种判断，即认为他的状态是“知道了”。}
    \explain{You drink too much.}{你喝酒喝太多了。}{这里用一般现在时，表达一种判断，即说话者主观认为“喝得太多了”是“你”的习惯。}
    \explain{He always finds out.}{他总会发现。}{这里用一般现在时，表达一种信念，即认为“他发现”是必然的事情。但实际上说话者可能并没有足够的观察，只是相信这个结论。}
\end{liju}

我们可以看到，语法既是一种规律，也是一种表达手段。利用语法规律，可以表达更具体的含义。这也是语法的功能。

\subsection{一般现在时的否定句}

否定一般现在时，就是说并没有这样的常识、真理、习惯、常见动作或存续状态。

如果只有普通动词，那么在主动词前加上助动词do，把主动词的变位转到do上，把主动词变回原形，最后在do后面加上not。
\begin{liju}
    \enzh{He teaches you English. \warr He does not teach you English.}{他教你英语。 \warr 他不教你英语。}
    \enzh{He knows. \warr He does not know.}{他知道了。 \warr 他不知道。}
    \enzh{We go shopping every weekend. \warr We do not go shopping every weekend.}{我们每周末都去购物。 \warr 我们不是每周末都去购物。}
    \enzh{He always finds out. \warr He does not always find out.}{他总会发现。 \warr 他不总会发现。}
    \enzh{All the customers drink tea. \warr All the customers do not drink tea.}{所有的顾客都喝茶。 \warr 不是所有的顾客都喝茶。}
\end{liju}

要注意的是：当主语有“所有”、“任何”这样涵盖全体的限定词（全肯定）时，用以上方式得到的否定句表达有否定。这在逻辑上是矛盾关系。而按照中文的习惯，形式上的否定会导致反对关系（全肯定和全否定）。

\subsection{一般现在时的疑问句}

前面举的例子都是陈述句。那么，一般现在时的疑问句是怎样的呢？

举例来说，我们想知道“老虎吃肉吗？”这种问题的答案。那么由于答案是“是”或者“不是”，要使用一般疑问句。我们这样提问：

\begin{liju}
    \enzh{Are they students?}{他们是学生吗？}
    \enzh{Is this true?}{这是真的吗？}
    \enzh{Can he swim?}{他会游泳吗？}
    \enzh{Shall I leave?}{我该离开吗？}
    \enzh{Do tigers eat meat?}{老虎吃肉吗？}
    \enzh{Does light travel faster than sound?}{光跑得比声音快吗？}
\end{liju}

主动词是be动词时，把be动词移到句首，其他成分保持不变。主动词不是be动词时，主动词位置不变。如果有情态动词，就把情态动词放在句首；如果没有情态动词，就在句首加Do或Does（取决于主语是否是第三人称单数），并把主动词改回原形。

对于特殊疑问句，我们只说针对谓语动词的特殊疑问句。对于普通动词，如果没有情态动词，根据一般规则，我们要把动词改成do，并做对应变位；然后在主语前加do，并把变位转移到它身上，最后在句首加what。例如：

\begin{liju}
    \enzh{He teaches English in Beijing. \warr He does. \warr Does he do. \warr What does he do?}{他在北京教英语。 \warr 他做什么？}
\end{liju}

如果有情态动词，那么把主动词改为do后，把情态动词移到主语前，再在句首加what。例如：

\begin{liju}
    \enzh{He can swim. \warr He can do. \warr Can he do. \warr What can he do?}{他会游泳。 \warr 他会做什么？}
    \enzh{He must finish the test. \warr He must do. \warr Must he do. \warr What must he do?}{他必须完成测试。 \warr 他必须做什么？}
    \compare{He can have a room. \warr What can he have?}{He can have a room. \warr What can he do?}
\end{liju}

\section{现在进行时}

上一节我们学习了一般现在时。一般现在时并不用来表示正在发生的动作。那么怎么表示正在发生的动作呢？我们可以用现在进行时。

现在进行时是什么样子的？我们以一般现在时为基准，把动词转换成现在进行时。具体做法是：首先在主动词前面加be动词，并把主动词的变位转到be动词上，然后再把主动词变为现在分词。我们先来看一个例子：

\vspace{8pt}
\begin{liju}
    \enzh{He studies English. \warr He is study English. \warr He is studying English.}{他学英语。 \warr 他正在学英语。}
\end{liju}

上面的句子里，我们从一般现在时出发，先在“studies”前面加上be，并把主动词的变位转到be上，然后再把主动词变为现在分词。这样就得到了现在进行时。

什么是现在分词呢？现在分词是动词的一种变形。由于变形后的一大作用是作为现在进行时的标记，所以叫现在分词。动词如何变为现在分词呢？规则是在词尾添加\phonetic{ɪŋ}。从拼写来说，就是在词尾添加\jie{ing}。如果动词以不发音的e结尾，那么去掉e，再在末尾添加\jie{ing}。如果动词最后一个闭音节是重读闭音节，且末尾恰有一个辅音字母，则需要双写该辅音字母，再在末尾添加\jie{ing}。

\begin{liju}
    \enzh{I study English. \warr I am studying English.}{我学英语。 \warr 我正在学英语。}
    \enzh{You play basketball. \warr You are playing basketball.}{你们在打篮球。 \warr 你们正在打篮球。}
    \enzh{She takes an exam. \warr She is taking an exam.}{她在考试。 \warr 她正在考试。}
    \enzh{He puts the book back on the shelf. \warr He is putting the book back on the shelf.}{他把书放回书架上。 \warr 他正把书放回书架上。}
\end{liju}

\subsection{现在进行时的否定句}

否定现在进行时，就是说并没有正在做某件事。具体的方式是在be和主动词之间加上not。
\begin{liju}
    \enzh{I am not teaching you English.}{我没在教你英语。}
    \enzh{You are not getting sick.}{你不会得病的。}
    \enzh{He is not watching TV.}{他并不在看电视。}
    \enzh{They are not laughing at the dog.}{他们并不在笑话那只狗。}
\end{liju}

\subsection{现在进行时的疑问句}
针对现在进行时的提问主要是一般疑问句。如一般规则所说，由于有助动词be，我们把be放在句首，其他词位置和顺序不变。

\begin{liju}
    \enzh{They are pushing the table. \warr Are they pushing the table?}{他们在推桌子。 \warr 他们在推桌子吗？}
    \enzh{He is telling you the truth. \warr Is he telling you the truth?}{他把真相说给你听。 \warr 他告诉你的是真相吗？}
    \enzh{The car is moving forward. \warr Is the car moving forward?}{车在往前移动。 \warr 车在往前移动吗？}
\end{liju}

对于针对谓语动词的特殊疑问句，根据一般规则，我们要把动词改成do，然后把助动词be移到主语前，最后在句首加what。例如：

\begin{liju}
    \enzh{He is telling you the truth. \warr He is doing. \warr Is he doing. \warr What is he doing?}{他把真相说给你听。 \warr 他在做什么？}
\end{liju}

\section{时态与时间}

从一般现在时和现在进行时的用法中可以看到，英语动词的时态和它们的名称不全相同，一般现在时往往并不是“一般”用来描述“现在”动作的时态，不能望文生义。要更好地理解时态，就要更细致地讨论语言里表达动作和时间关系的方法。

中文里并没有时态的概念。我们说中文时，往往只在需要的时候，用时间词来说明动作发生的时间。如果动作和时间的关系不紧密、不重要，我们就不一定强调它。比如：我们说“陈老师在立德中学教英语”，也说“陈老师教我句首字母要大写”。从英语的思维来看，两句话里“教”的时态不同。但在中文里，我们并不作区分，因为“教”和时间的关系在句子里并不是最重要的信息。如果我们要强调“教”和时间的关系，可以说“陈老师现在在立德中学教英语”，也说“陈老师昨天教我句首字母要大写”。

英语，和很多欧洲语言一样，选择用动词变位（动词形态的屈折变化）来表达动作和时间的关系。而由于变位是语法规律，所以所有动词在做主动词的时候，都要变位（或者由情态动词、助动词进行变位），确定一种时态。因此，英语就必须考虑很多细致的具体问题，为很多中文里不必要区分的情况做出规定，该用什么时态。有些时候，这种规定仅仅是因为必须给动词发一顶时态的“帽子”，所以并不一定完全合理，只是反映大众的使用习惯。它只是我们解释大众的使用习惯时给出的“道理”。

一种常见的“道理”是所谓的“动作发生时间”和“叙述的时间”。一个句子里的动作，根据句子内容含义，有一个发生的时间。而我们写出、讲出这个句子，也有一个时间。这两个时间是不一样的。英语语法的时态解释里，以后者为“现在”。所以，如果动作发生时间不是“叙述的现在”，那么就不是“现在时”。比如，假设我们是负责盯梢的警察，我们在警车里看到一个人走进一栋建筑里，可以用中文描述说：一个穿黑衣服的男人走进酒店，一个女人在大堂里等他。这句话里，不需要强调“走”和“等”的时间。这是中文的习惯。在英语里，我们说“A man in black walked into the hotel. A woman is waiting for him in the lobby.”这里的动作“walk”和“wait”都需要加上时态，并且不能错，否则会改变句子的含义。在这句话里，动作“walk”发生于叙述前，所以用一般过去时，而叙述时“wait”仍然在进行，所以用现在进行时。

有了“动作发生时间”和“叙述的时间”这两个概念，理解英语的时态就容易多了。



\section{一般将来时}

如果说一般现在时并不用来表示现在的动作，那么一般将来时倒确实用来表示将来要发生的动作或预计会发生的动作。它的形态是“will+动词原形”。这个形式对任何人称和单复数的主语都适用。

\begin{liju}
    \enzh{I will teach you English.}{我将教你英语。}
    \enzh{You will be happy.}{你会快乐的。}
    \enzh{He will give you the report tomorrow.}{他明天会把报告给你。}
    \enzh{They will shut down this computer.}{他们会把这台电脑关掉。}
\end{liju}

要注意的是：
\begin{enumerate}
    \item be动词的原形是be，因此一般将来时是will be。
    \item 一般将来时的句子没有情态动词，有情态动词的句子没有一般将来时。要表达将来的态度、意愿、情势等等，需要用其他的方式。详见以后的介绍。
\end{enumerate}

\subsection{一般将来时的否定句}

否定一般将来时，就是说将来不会发生某件事。具体的方式是在will和主动词之间加上not。
\begin{liju}
    \enzh{I will not teach you English.}{我不会教你英语。}
    \enzh{You will not be happy.}{你不会快乐的。}
    \enzh{He will not give you the report tomorrow.}{他明天不会把报告给你。}
    \enzh{They will not shut down this computer.}{他们不会把这台电脑关掉。}
\end{liju}

\subsection{一般将来时的疑问句}

针对一般将来时的提问主要是一般疑问句。如一般规则所说，由于有助动词will，我们把will放在句首，其他词位置和顺序不变。

\begin{liju}
    \enzh{They will get the prize. \warr Will they get the prize?}{他们会获奖。 \warr 他们会获奖吗？}
    \enzh{He will give you a book. \warr Will he give you a book?}{他会给你一本书。 \warr 他会给你一本书吗？}
    \enzh{The car will hit the tree. \warr Will the car hit the tree?}{车要撞上树了。 \warr 车要撞上树了吗？}
\end{liju}

对于针对谓语动词的特殊疑问句，根据一般规则，我们要把动词改成do，然后把助动词will移到主语前，最后在句首加what。例如：

\begin{liju}
    \enzh{He will give you a book. \warr He will do. \warr will he do. \warr What will he do?}{他在北京教英语。 \warr 他做什么？}
\end{liju}

\section{一般过去时}

一般现在时并不用来表示具体发生的动作，但我们可以用一般过去时来做到这一点。

一般过去时的形态和动词原形不一样，一般来说要在词尾加上\phonetic{d}、\phonetic{t}或\phonetic{ɪd}音表达区别。具体规则如下：

\begin{enumerate}
    \item 动词以\phonetic{p}，\phonetic{k}，\phonetic{f}，\phonetic{θ}、\phonetic{s}，\phonetic{ʃ}，\phonetic{tʃ}等辅音结尾则后加\phonetic{t}，
    \item 动词以\phonetic{b}，\phonetic{g}，\phonetic{n}，\phonetic{l}，\phonetic{m}，\phonetic{n}，\phonetic{ŋ}，\phonetic{ð}，\phonetic{r}，\phonetic{z}，\phonetic{ʒ}，\phonetic{dʒ}等辅音以及元音结尾则加\phonetic{d}，
    \item 动词以\phonetic{t}、\phonetic{d}结尾则加\phonetic{ɪd}，
    \item 不规则动词：单列说明。
\end{enumerate}

拼写上表现为：
\begin{enumerate}
    \item 大多数情况下，在词末直接加上\jie{ed}，表示\phonetic{d}、\phonetic{t}或\phonetic{ɪd}音；
    \item 以不发音的e结尾的词则只加\jie{d}；
    \item “辅音字母+\;y”结尾的动词，改y为i再加\jie{ed}，\jie{ied}读\phonetic{ɪd}；
    \item “元音字母+\;y”结尾的动词，直接加\jie{ed}；
    \item 如果动词最后一个音节是重读闭音节，且末尾只有一个辅音字母，则需要双写该辅音字母，然后再加\jie{ed}；
    \item 不规则动词：单列说明。
\end{enumerate}

一般过去时变位不规则的动词较多，详见附录中的《不规则动词变位表》。

什么场合下用一般过去时呢？对于在过去发生的动作、事情，过去的状态，要用一般过去时。它强调的是过去已经发生的事情，而不是对现在的影响。 

\begin{liju}
    \explain{I saw her in the street yesterday.}{昨天我在街上看见她了。}{过去某个时间点发生的事情。}
    \explain{She was a lawyer.}{她以前是律师。}{过去的状态。}
    \explain{We went shopping every weekend.}{我们以前每周末都去购物。}{过去经常做的事，过去的习惯。}
\end{liju}

当我们使用一般过去时描述过去的经常动作、习惯时，往往暗示着现在并不这样了，但也并不表示确定的否定。

要注意的是，使用情态动词表示意愿、态度、能力、情势的句子，不一定能通过语法上的变位来表达“过去的意愿、态度、能力、情势”。这时候，有些情态动词要转为对应含义的普通动词。

\subsection{一般过去时的否定句}

否定一般过去时，就是说过去没有发生某事、没有某状态、没有某经常动作或习惯。一般过去时变位没有助动词，所以否定的方式和一般现在时同理。
\begin{liju}
    \enzh{I saw her in the street yesterday. \warr I did not see her in the street yesterday.}{昨天我在街上看见她了。}
    \enzh{She was a lawyer. \warr She was not a lawyer.}{她以前是律师。 \warr 她以前不是律师。}
    \enzh{She could do the splits when she was a kid. \warr She could not do the splits when she was a kid.}{她小时候能劈叉。 \warr 她小时候不能劈叉。}
    \enzh{He might be at home last night. \warr He might not be at home last night.}{他昨晚也许在家。 \warr 他昨晚也许不在家。}
\end{liju}

\subsection{一般过去时的疑问句}

针对一般过去时的提问主要是一般疑问句。一般过去时变位没有助动词，所以变化方法和一般现在时同理。

\begin{liju}
    \enzh{They took the test yesterday. \warr They did take the test yesterday. \warr Did they take the test yesterday?}{他们昨天做了测试。 \warr 他们昨天做了测试吗？}
    \enzh{She was a lawyer. \warr Was she a lawyer?}{她以前是律师。 \warr 她以前是律师吗？}
    \enzh{She could do the splits when she was a kid. \warr Could she do the splits when she was a kid?}{她小时候能劈叉。 \warr 她小时候能劈叉吗？}
\end{liju}

对于针对谓语动词的特殊疑问句，变化方法和一般现在时同理。例如：

\begin{liju}
    \enzh{He gave you a book. \warr He did. \warr  Did He do. \warr  What did he do?}{他给了你一本书。 \warr 他做了什么？}
    \enzh{She could do the splits when she was a kid. \warr She could do \warr  Could she do. \warr  What could she do?}{她小时候能劈叉。 \warr 她以前能做什么？}
\end{liju}

\chapter{句子与单词之二}

\section{词汇之旅——英语词汇的来源}

请先观察以下几组单词：
\begin{center}
    \begin{minipage}{0.6\linewidth}
    \centering
        \begin{itemize}
            \item ask（问）和question（质问）
            \item fast（快）和rapid（迅速）
            \item get（得）和obtain（获得）
            \item sea（海）和ocean（海洋）
        \end{itemize}
    \end{minipage}
\end{center}

我们会发现，每组第一个单词更简单短小，含义更基本；第二个单词音节较多，含义更具体。

再看一组例子：
\begin{center}
    \begin{minipage}{0.6\linewidth}
    \centering
        \begin{itemize}
            \item 家畜：pig\phantom{k}（猪），\phantom{肉}cow（牛），\phantom{肉}sheep\phantom{jj}（羊）
            \item 肉类：pork（猪肉），beef（牛肉），mutton（羊肉）
        \end{itemize}
    \end{minipage}
\end{center}

我们会发现：中文里“猪”和“猪肉”的关系一目了然，英语里两个词却长得完全不一样。

这是因为在漫长而曲折的发展历程中，英语吸收了来自不同文化和语言的词汇。我们可以将英语词汇的积累看作一场跨越千年的旅程。它先后经历了三个重要的历史阶段，形成了今天丰富多元的面貌。

第一个阶段大约从公元5世纪开始，盎格鲁–撒克逊人（the Anglo-Saxons）将他们的语言带到了不列颠群岛，也就是如今的英国。盎格鲁–撒克逊人是古代日耳曼人（the Germanic people）的一支，他们说一种日耳曼语。盎格鲁–撒克逊人的语言构成了英语最古老、最核心的词汇（我们称作日耳曼词）。它们往往只有一到两个音节，简单短小，主要包括（举例）：
\begin{itemize}
    \item 自然万物：water（水），fire（火），sun（日），stone（石头），tree（树），sea（海）
    \item 身体部位：hand（手），eye（目），ear（耳），heart（心），head（头），foot（脚）
    \item 日常生活：home（家），house（屋），door（门），ground（地），child（童），food（食），bed（床），road（路），cloth（布），pig（猪），cow（牛），sheep（羊），cat（猫），mouse（鼠）
    \item 基本动作：go（去），come（来），see（看），make（作），eat（吃），drink（喝），ask（问）
    \item 基本修饰语：fast（快），red（红），tall（高），good（好），hard（硬），sad（悲），dark（暗）
    \item 最基本的语法词：the，a，I，us，he，she，this，that等等
\end{itemize}

这些词汇是英语的“骨架”，是日常生活沟通必需懂得的词汇。它们和德语以及一些欧洲北方民族的词汇有相似之处，因为它们都传自古代日耳曼人。

1066年，诺曼人（the Normans）征服了不列颠群岛。这是英语发展的转折点。诺曼人是今天法国人的祖先之一。其后的数百年里，古法语成为英国的官方语言。在前面的例子里，家畜用日耳曼词，而上餐桌的肉食则是法语词。这是因为英国农民饲养家畜，而统治他们的诺曼贵族享用餐桌上的美食。

法语的词汇又多来自拉丁语（Latin）。因此，大量拉丁语词经由法语传入英语，也有很多通过学术、宗教等途径直接进入英语。法语和拉丁语词极大地丰富了英语的词汇，主要包括（举例）：
\begin{itemize}
    \item 社会、政治与法律：government（政府），court（法庭），justice（正义），official（官员）
    \item 都市、文明与艺术：manners（礼节），restaurant（餐厅），art（艺术），market（市场）
    \item 思想、学术和教育：university（大学），science（科学），logic（逻辑），lesson（课程）
    \item 抽象概念：message（信息），success（成功），dialog（对话），color（颜色），patience（耐心）
    \item 细致和抽象的动作：prefer（偏好），receive（接收），purchase（购买），explain（解释）
    \item 细致和抽象的修饰语：vivid（栩栩如生），honest（诚实），elegant（精致），divine（神圣）
\end{itemize}

这一层次的词汇通常音节更多，很多是正式的书面语，它们赋予了英语表达抽象概念和复杂思想的能力。

要说明的是，日耳曼语、法语、拉丁语和英语本身也在不断演变。以上的单词源自古代的日耳曼语、法语、拉丁语，但进入英语时往往不是今天我们看到的样子。事实上，我们今天学习的英语（也叫现代英语）直到16世纪才逐渐成型。而我们现在看到的这些单词是演变后的样子。

16世纪以后，现代英语逐渐成型，而随着不列颠帝国的全球殖民扩张，英语走向世界，同时也像海绵一样，从各大洲的殖民地语言中吸收了大量新词。不仅如此，英语更是在全球化进程中持续从各国语言中借用词汇。比如：
\begin{itemize}
    \item 来自美洲语言：canoe（独木舟）、jerky（肉干）、tobacco（烟草）
    \item 来自印度语：shampoo（洗发水），bungalow（平房），pyjamas（睡衣）
    \item 来自汉语：tea（茶），typhoon（台风）、kowtow（叩头）
    \item 来自日语：tsunami（海啸）、ninja（忍者）、origami（折纸）
\end{itemize}

由于这个时期的英语已经逐渐成熟，我们把后来进入的词称为“借词”，即从外语借来的词。它们展现了英语文化与世界各地区文化交融的现状。经过数百年的演变，英语已经不是不列颠群岛或欧洲的语言，而是世界性的语言。英语影响了众多语言，也受到众多语言的影响。

\section{英语的声音密码——节奏、连读与语调}

听下面两个句子，感受它们的不同：
\begin{center}
    \begin{minipage}{0.6\linewidth}
    \centering
        \begin{itemize}
        \item[中文：] 各位先生，各位女士，晚上好！
        \item[英文：] Good evening, ladies and gentlemen!
        \end{itemize}
    \end{minipage}
\end{center}

你是否觉得，中文的每个字发音都独立、饱满，而英文的句子听起来却像波浪一样有起伏、有快慢，连绵不绝？这种差异，正是我们今天要探讨的“英语的声音密码”。掌握这些密码，能让你更容易听懂英语，说起英语来也更自然。

我们知道，汉语的单位是汉字。每个汉字都是独立的方块字，对应一个音节。我们说汉语的时候，就是把一个个音节清晰完整地发出来。说汉语的时候，我们讲究字正腔圆。这里的“圆”就表示每个汉字对应的音节都是一个独立的单元，我们把这个音节发得圆满了，就发好了音。

英语和汉语不一样。英语是以音为本位的语言。音节连接成单词，单词又连接成句子。虽然在书面记录时，我们用空格来分隔单词，但说话的时候，把单词分隔开来发音，就不地道了。英语的单词里往往有多个音节，因此要说得很快。要是刻意把单词一个一个来发音，中间用停顿分隔，就来不及了。

说英语的人，习惯把单词和单词连起来读，再加上重音和起伏来标记节奏。我们可以形象地把汉语朗读想象成一列大小不一的圆，各自分开；而英语朗读就好比一串波浪般的圈圈，大小不一，每个圈和前后的圈首尾相连。用书法来比拟，汉语朗读就好比正楷，每个字都方方正正，各安其位，各不相及；英语朗读就好比行草，讲究笔韵流转，气息不绝，藕断丝连。

\subsubsection*{节奏的源泉：重读和弱读}

在不止一个音节的英语单词中，总有一个音节读得比其他音节更响亮、更饱满、音调也稍高，这就是单词的重读音节。它就像一个单词的“心跳”。下面我们用大写字母来表示重读音节，用小写字母表示非重读音节。

\begin{itemize}
    \item \textbf{TEA}cher（教师）
    \item \textbf{STU}dent（学生）
    \item \textbf{AP}ple（苹果）
    \item \textbf{BEAU}tiful（美丽的）
\end{itemize}

如果你把重音放错了位置，比如把 APple 读成 apPLE ，听起来就会很别扭，甚至造成误解。这与中文的“苹果”两个字音量相当的特点完全不同。

单词组成句子，句子也有自己的节奏。在句子中，那些传递核心信息的词通常需要重读，读得缓慢清晰；而那些不那么重要或者只起语法作用的词则会弱读，读得又轻又快。例如：
\begin{itemize}
    \item The \textbf{TEA}cher has a \textbf{BOOK}.
    \item \textbf{I} \textbf{LOVE} my \textbf{FAM}ily.
\end{itemize}

这种分轻重缓急的模式，构成了英语的基本节奏，就像音乐中的强拍和弱拍交替出现。在音标中，我们会用“\phonetic{'}”放在重读的音节前，标记重音。非重音的音节就是弱读音节。

另外，在一个句子里，弱读的音节可能发生变化，和在单词里的读法不一样。比如，冠词“a”单独读\phonetic{eɪ}，但在句子里作为弱读音节时，会读成\phonetic{ə}。冠词“an”单独读\phonetic{æn}，但在句子里作为弱读音节时，会读成\phonetic{ən}。具体的变化方式，我们将在后面详细介绍。动词“is”单独读\phonetic{ɪs}，但在句子里作为弱读音节时，会读成\phonetic{ɪz}。具体的变化方式，我们将在后面详细介绍。

\subsubsection*{流利的秘诀：连读和缩读}

只靠轻重，还是无法形成波浪的。为了说话更省力、更流畅，快速说话时，英语单词的读音会发生变化。这是英语的习惯。变化主要有两种：连读和缩读。

首先来说连读。当相邻的两个单词，前一个以辅音结尾，后一个以元音开头时，人们通常会自然地将它们连在一起读。

以下是几个典型的例子\footnote{注意：箭头后的写法只用来示意读音如何改变，连读时不能按此改变写法。}：
\begin{itemize}
\item an apple（\phonetic{æn 'æpəl}） $\rightarrow$ a-\textbf{napple}（\phonetic{ə 'næpəl}）
\item my name is （\phonetic{maɪ 'neɪm ɪs}）$\rightarrow$ my na-\textbf{mis}（\phonetic{maɪ 'neɪmɪz}）
\item It is a dog. （\phonetic{'ɪt ɪs eɪ 'dɒg}）$\rightarrow$ I-\textbf{ti}-\textbf{sa} dog.（\phonetic{'ɪtɪzə 'dɒg}）
\end{itemize}

这不是一种强制的规则，而是长久以来，用英语说话的人为了说起来更快、更流畅而自然形成的习惯。

其次，在日常口语中，一些非常高频的词常常会和前面的词合并，形成缩读形式。这同样是语速加快的自然产物。

以下是几个典型的例子\footnote{注意：箭头后的写法是合乎语法的标准写法，缩读时必须按此改变写法。}：
\begin{itemize}
\item I am （\phonetic{'aɪ æm}）$\rightarrow$ \textbf{I'm}（\phonetic{'aɪm}）
\item It is （\phonetic{'ɪt ɪs}） $\rightarrow$ \textbf{It's}（\phonetic{'ɪts}）
\item they are （\phonetic{'ðeɪ ɑːr}）$\rightarrow$ \textbf{they're}（\phonetic{'ðeɪr}）
\item do not （\phonetic{'duː nɒt}）$\rightarrow$ \textbf{don't}（\phonetic{'dɒnt}）
\end{itemize}

请记住，缩读是口语中非常普遍的现象，是地道的英语，而不是错误。缩读也是约定俗成的，不是任何两个词都可以缩读的。

\subsubsection*{情绪的表达：语调起伏}

除此以外，语调，即说话时声音高低的变化，也很重要。“波浪”如何“起伏”，是表达情绪和态度的关键。同样一句话，用不同的语调说出来，意思可能完全不同。

降调，指声音从高到低下降。换句话说，就是全句最后一个“波浪”以“波谷”结束，通常表示肯定和结束。常用于：
\begin{itemize}
\item \textbf{陈述句}：表示肯定、陈述事实。\ My name is Tom. $\searrow$
\item \textbf{特殊疑问句}：询问具体信息。\ What's your name? $\searrow$
\item \textbf{祈使句}：表示命令。\ Sit down. $\searrow$
\end{itemize}

降调给人一种“话已说完，意思明确”的感觉。

升调，指声音从低到高上升。换句话说，就是全句最后一个“波浪”以“波峰”结束，表示疑问与不确定。常用于：
\begin{itemize}
\item \textbf{一般疑问句}：表示询问，期待对方回答“是”或“否”。\ Is this your book? $\nearrow$
\item \textbf{表示礼貌或不确定}：即便不是疑问句，用升调也可以让语气更委婉。\ Really? $\nearrow$（表示惊讶、怀疑。注意：降调的 Really? $\searrow$ 有“不以为然”的暗示。）
\end{itemize}

升调给人一种“话未说完，尚有疑问”的感觉。

英语的声音系统（重音、连读、语调）是一个有机的整体，它们共同作用，赋予了英语独特的节奏感和丰富的表现力。

\section{问候与礼貌：日常用语背后的文化心理}

阅读下面三个场景中的对话：

\textbf{场景一：}2018年，美国某个小镇的咖啡馆里，你遇到了与你约好的朋友Alex。
\begin{center}
    \begin{minipage}{0.6\linewidth}
    \centering
        \begin{itemize}
            \item[你：] Hey Alex! What's up? 
            \item[Alex：] Hey! Not much. Just grabbing my coffee. Good to see you!
        \end{itemize}
    \end{minipage}
\end{center}


\textbf{场景二：}1992年，英国某大学的校园里，你偶然遇到了正在散步的Bennett教授。
\begin{center}
    \begin{minipage}{0.6\linewidth}
    \centering
        \begin{itemize}
            \item[你：] Good afternoon, Professor Bennett.
            \item[Bennett教授：] Ah, good afternoon. Lovely day, isn't it?
            \item[你：] Yes, it is. How are you?
            \item[Bennett教授：] Very well, thank you.
        \end{itemize}
    \end{minipage}
\end{center}


\textbf{场景三：}1880年，英国伦敦某个沙龙里，Winchester女士遇到了她的朋友Ashworth先生。
\begin{center}
    \begin{minipage}{0.6\linewidth}
    \centering
        \begin{itemize}
            \item[Ashworth先生：] Lady Winchester. How do you do?
            \item[Winchester女士：] Mr. Ashworth. How do you do? It is a pleasure to see you.
        \end{itemize}
    \end{minipage}
\end{center}

仔细观察这三段对话，你会发现英语中的“打招呼”远不止是“说你好”那么简单。不同的场合、不同的人物关系，使用的问候语也完全不同。今天，我们就来破解英语问候语背后的文化密码。

问候语是什么？它就像是进入社交世界的一串“通用密码”。当你对别人说出问候语时，你实际上是在发送三个重要的信号：“我看到了你。我确认你的社会地位。我没有敌意。”

这个信号在不同的文化里，用不同的方式发送。比如在中文环境里，我们通常不会对陌生人主动问候，问候更多发生在熟人之间。而在英语文化中，即使是陌生人之间，在电梯里、在小路上相遇，也常常会微笑点头或者说一声“Hi”或“Good morning”。这是一种普遍的社交礼仪。正因如此，学习英语的问候方式对我们来说才如此重要。

从上面的场景中，我们可以提炼出理解英语问候文化的三把“钥匙”。

\subsubsection*{关系的亲疏}

关系越亲近，问候就越简单随意；关系越正式疏远，问候就越复杂客气。
\begin{itemize}
\item \textbf{好友之间}（场景一）：用“Hey”“What's up?”这样简短的词，就像我们说的“嗨”“干嘛呢”。
\item \textbf{师生之间}（场景二）：用“Good afternoon”“How are you?”这样完整的句子，并一定要使用头衔“Professor”。
\item \textbf{正式场合}（场景三）：用“How do you do?”这样现在已经很少用的固定句式，并且伴随鞠躬等肢体语言。
\end{itemize}

要注意的是：在英语中，对不熟悉的人直呼其名是不礼貌的，通常要加上Mr., Ms.等称呼。如果对方有头衔，应当使用头衔进行称呼，表示尊重。

\subsubsection*{场合的正式程度}

场合越正式，用语就越固定，就像在念一个“标准剧本”。
\begin{itemize}
\item \textbf{随意场合}：可以说“Hi”、“Hello”、“How's it going?”区别不大，对方不会在意。
\item \textbf{正式场合}：要说“Good morning\,/\,afternoon”、“How do you do?”（过时）。对方会根据你的选择来判断你的意图和态度，并作出回应。
\item \textbf{万能问候}：如果不确定该用什么的话，最好用比较正式、客气的方式打招呼，以免产生误会。“How are you?”在大多数非正式和半正式场合都可以使用，是比较安全的打招呼方式。在不正式的场合或年轻人之间，对方可能觉得你有些滑稽古板。但如果在较正式的场合使用“How's it going?”、“What's up?”这样的说法，会给人留下不佳的印象。
\end{itemize}

\subsubsection*{文化的习惯：积极的回应}

这是英语问候文化中非常特别的一点。请注意看场景一和场景二的回答：
\begin{center}
    \begin{minipage}{0.6\linewidth}
    \centering
        \begin{itemize}
            \item 问：What's up? （有什么新鲜事吗？）\\
            答：Not much.（没啥。）
            \item 问：How are you? （你好吗？）\\
            答：Very well, thank you. （很好，谢谢。）
        \end{itemize}
    \end{minipage}
\end{center}

你会发现，无论实际情况如何，人们通常都会给出一个简短、中性或积极的回答。这并非虚伪客套，而是英语问候文化中的一个重要习惯。这个习惯的背后，是对“个人边界”的尊重——不过多用自己的琐事或烦恼去打扰对方，让社交氛围保持轻松和积极。因此，当被问到“How are you?”时，回答“Good.”或“Fine, thanks.”、“Not bad.”是最常见、最安全的选择。

这一点上，可以看出英语文化与中国文化的差异。中国文化中，关心他人或相互交流私密信息，是释放善意的一种方法。但在英语文化中，向他人打听，或透露私密信息，反而是一种冒犯。这种差异，会体现在在打招呼、问候，乃至一般礼仪上。比如，在场景二里，我们可以看到，英国人相互问候时喜欢聊天气，因为天气是一种不涉及个人私密信息的话题。

除了问候语，一些常用的“小词”同样重要：
\begin{itemize}
\item \textbf{Please}：这是表达“请求”而非“命令”的关键词。忘记使用Please，即使语法正确，也会显得很粗鲁。
\item \textbf{Thanks\;/\;Thank you}：使用频率非常高，无论多小的帮助都值得一声“Thank you”。
\item \textbf{Sorry\;/\;Excuse me}：两者有明确分工。\textbf{Excuse me}用于事前打扰（借过、提问前），\textbf{Sorry}用于事后道歉（不小心碰到别人）。
\end{itemize}

这些词就像机器里的润滑油，能让社会交往更顺畅。除了这些简单的“小词”，英语里还有很多更复杂的方法来表达善意和礼貌。我们将在后面详细介绍。

语言是文化的镜子。一种语言的表达方式，反映了使用这种语言的人们的思维方式和社会规则。我们的目标不仅仅是说出语法正确的句子，更是要在合适的场合，对合适的人，说出得体的话。刚开始时，你可能会觉得不习惯，但只要你带着真诚和善意，就不妨勇敢地说出来。


\section{定语和状语}

\section{形容词和副词}

\section{归属的代词}

表示归属的方法：

\begin{liju}
    \enzh{This is a pen.}{这是一支笔。}
    \enzh{This is my pen.}{这是我的笔。}
    \enzh{This is your pen.}{这是你的笔。}
    \enzh{This is our pen.}{这是我们的笔。}
    \enzh{This is your pen.}{这是你们的笔。}
    \enzh{This is his pen.}{这是他的笔。}
    \enzh{This is her pen.}{这是她的笔。}
    \enzh{This is their pen.}{这是他们的笔。}
    \enzh{This is its pen.}{这是它的笔。}
\end{liju}

另外，“这个/这些”和“那个/那些”也有归属形式：

\begin{liju}
    \enzh{This pen is red.}{这支笔是红笔。}
    \enzh{These pens are red.}{这些笔是红笔。}
    \enzh{That pen is blue.}{那支笔是红笔。}
    \enzh{Those pens are blue.}{那些笔是红笔。}
\end{liju}

归属代词指代形容词。

直接表示归属物的方法：

\begin{liju}
    \enzh{This pen is mine.}{这支笔是我的。}
    \enzh{This pen is yours.}{这支笔是你的。}
    \enzh{This pen is ours.}{这支笔是我们的。}
    \enzh{This pen is yours.}{这支笔是你们的。}
    \enzh{This pen is his.}{这支笔是他的。}
    \enzh{This pen is hers.}{这支笔是她的。}
    \enzh{This pen is theirs.}{这支笔是他们的。}
    \enzh{This pen is its.}{这支笔是它的。}
\end{liju}

属物代词指代名词。

\section{介词和连词}

\section{数词}

\section{英语拼写之谜：声音与文字之间的历史鸿沟}

试读一下这个单词：\textbf{knight}。它读作 \phonetic{naɪt}。为什么它的拼写里会有 “k” 和 “gh” 这些不发音的字母呢？我们说，英语的拼写是对读音的记录，但很多英语单词的发音和拼写都“对不上”。这是为什么呢？下面我们就来解开这个谜团。

英语拼写不是一张精确的“发音地图”，而是一座历史的博物馆。许多看似奇怪的拼写，其实都是历史的见证。从英语拼写的历史来看，有三个因素塑造了今天英语拼写的样貌。

\subsubsection*{元音大推移}

长久以来，英语的拼写方式不是统一的。古时候，各地乃至各个文人都随意地使用自己喜欢的拼写方式来记录英语的发音。大约从15世纪初到17世纪末，英语发生了一次重大的语音变化，语言学家称之为“元音大推移”（the Grand Vowel Shift）。简单来说，就是绝大多数长元音的发音方式都发生了改变。

几乎在同一时期，印刷术从欧洲大陆传入英国。这是英语发展的关键转折点。印刷术让书面的英语记录数量急剧增长。印刷一版书籍，就代表数千乃至数万份书面复制，比以往用手抄录快得多。手写的书籍往往只在私下流转，每次抄录都形成独特的版本。而印刷品可以轻易流传到千百里外，让千万人同时看到完全一样的内容。所以，印刷品的传播能力和影响力远超手写。而印刷时决定的单词拼写方式，就会对大众产生巨大影响。于是，许多单词的拼写，就在其发音即将发生巨变的前夕，被“冻结”在了印刷品上。

让我们回到开头的例子：\textbf{knight}。14世纪前，它的发音接近 \phonetic{kniht}，其中的 “k” 和 “gh” 都是发音的。印刷术冻结了 “k-n-i-g-h-t” 这个拼写。随后几百年里，它的发音持续简化：“k” 不读了，“gh” 的音也消失了，元音也变了，最终成了今天的发音： \phonetic{naɪt}。所以，knight 这个拼写就像一块语言化石，记录的是它在14世纪以前的发音。

\subsubsection*{混搭的规则}

英语是一座著名的“词汇大熔炉”。它从法语、拉丁语、希腊语等语言中借用了大量词汇。借词时，英语常常把单词的拼写和意义一起借来，但这些拼写反映的是外语的发音规则，而不是英语的发音规则。说英语的人试图用英语的方式来读这些单词，但外语的拼写往往不能完全用英语规则读出。有些能用英语发音的外语规则进入了英语，成为英语发音规则的一部分。这样，同样的拼写，有些词就要用外语的发音规则来读了。但如果外语的发音比较困难，说英语的人就会用相近的英语发音来代替。这种种做法，导致英语里出现了“一套字母，多种规则”。多种发音规则混杂在一起。

例如：
\begin{itemize}
    \item 来自法语的 \textbf{machine}（机器），保留了法语拼写和发音习惯：\jie{ine}结尾读\phonetic{i:n}成为了英语的一种发音规则。但它与更古老的发音规则冲突了。所以我们会看到 line（直线）读\phonetic{laɪn}，而 machine，routine（常规），cuisine（美食），marine（海洋的）这样的法语词里读\jie{ine}结尾读\phonetic{i:n}。
    \item 来自希腊语的 \textbf{psychology}（心理学），开头的 “p” 在希腊语中发音，英语借来时也保留了拼写，但“ps”音开头在英语中罕见，因此英语抛弃了“p”音。
    \item 来自法语的 \textbf{voyage}（航行），保留了法语拼写，并试图模仿法语的发音。法语原本读\phonetic{vwajaʒ}，但英语没有对应的读音，因此改为读\phonetic{ˈvɔɪɪdʒ}。
\end{itemize}

不同来源的词汇带来了各自语言的拼写和发音规则，它们共用一套字母，规则“打架”就在所难免了。

\subsubsection*{规范改革难}

18世纪以前，英语拼写相当随意，同一个词可以有多种写法。18世纪时，随着启蒙思想的传播，“百科全书派”学者开始系统地整理并规范各种知识。塞缪尔·约翰逊等学者开始编纂英语的辞典，希望系统地规范英语拼写。问题是，依据什么来选择一个标准拼法呢？他们遇到了一个大问题：英语世界从未建立统一的权威机构来管理语言。与法语有法兰西学术院、西班牙语有皇家学院不同，英语的标准化主要依赖学者个人的努力。因此，他们采取的标准是混杂的：
\begin{itemize}
    \item \textbf{词源原则}：根据单词的来源确定拼写。但由于能力和资料不足，颇多贻误。比如 island（岛屿）中不发音的 “s”，就是当时学者们认为它源于拉丁语 “insula” 而硬加进去的。后来证明这是个错误。
    \item \textbf{历史原则}：保留历史上流传下来的传统拼法。比如 psychology（心理学）开头的“p”已经不读，但仍旧保留。
    \item \textbf{发音原则}：在部分词中，让拼写尽量贴近当时的实际发音。比如 order（命令）来自法语 ordre，但原本法语的发音\phonetic{ˈɔʁdʁə}对说英语的人来说比较困难，所以在英语中，ordre的发音逐渐变为\phonetic{ˈɔrdər}。后来，拼写也转为和新发音更一致的 order。类似的单词还有tender（源自tendre）、member（源自membre）等。
\end{itemize}

更复杂的是，随着不列颠帝国在全球殖民扩张和美国独立，英语在不同地区各自独立发展，没有权威的机构来规范全球各地英语的书写。因此，任何统一的英语拼写改革都不可能发生。我们今天学习的英语拼写，是历史上人为选择与自然演化后多方妥协的结果，虽然大致取得了统一，但历史层积下来的种种问题已经无法消除了。

了解了英语拼写的历史，我们学习英语单词的拼写就不会摸不着头脑了。遇到拼写和发音对不上的单词，我们就能猜到，它要么来自古老的记录，要么来自英语以外的语言。随着掌握的英语单词不断增加，我们还可以对这些单词分门别类，在它们身上重新找出历史的痕迹，总结出一些局部的规律，方便记忆。比如，我们可以发现，某些拼写有相同的历史渊源，因此发音相同，可以看作一个“家族”：
\begin{itemize}
    \item \jie{ight} \textbf{家族}：knight，night，light，fight，right，sight，might，tight。它们都发 \phonetic{aɪt} 的音。
    \item \jie{tion} \textbf{家族}：action，nation，notion，motion，function，station，它们都发 \phonetic{ʃən}音。
    \item \jie{age} \textbf{家族}：baggage，message，village，average，courage，passage，它们都发 \phonetic{ɪdʒ}音。
    \item \jie{ceive} \textbf{家族}：receive，perceive，deceive，conceive，它们都发 \phonetic{si:v}音。
\end{itemize}

又比如，一个词进入英语后，虽然被英语发音规则改变了读音，但拼写和意义仍然保留下来。我们可以“顺藤摸瓜”，根据拼写的残留，找出含义的变迁和演化，帮助记忆。比如，sign（徵）来自法语，它在英语中不断变化，衍生出signal（信号）、signature（签名）、significant（显著的）、design（设计）、assign（分配）、resign（辞职）等词语。这些词语中“sign”的读音虽然不完全相同，但它们的含义都和sign有关。

那么，英语拼写里具体有哪些局部规律呢？我们会在以后详细介绍。

\chapter{}

\end{document}
